% Sintesi del paper e raccolta degli enunciati.
% Compilare da paper/summary/ dopo aver compilato ../main.tex (serve ../main.aux):
%   python3 make_statements.py && latexmk -pdf summary.tex
\documentclass[11pt]{article}
\usepackage[utf8]{inputenc}
\usepackage[T1]{fontenc}
\usepackage{amsmath,amssymb,amsthm,mathtools}
\usepackage{bm}
\usepackage[margin=1in]{geometry}
\usepackage[expansion=false]{microtype}
\usepackage{enumitem}
\usepackage{xr}
\externaldocument{../main}
\usepackage[capitalise]{cleveref}

\crefname{hypothesis}{Hypothesis}{Hypotheses}
\Crefname{hypothesis}{Hypothesis}{Hypotheses}
\crefname{conjecture}{Conjecture}{Conjectures}
\Crefname{conjecture}{Conjecture}{Conjectures}

\providecommand{\texorpdfstring}[2]{#1}  % hyperref is not loaded
% Shared macros for the DFL paper. Do not redefine these in section files.
\newcommand{\E}{\mathbb{E}}
\newcommand{\Prob}{\mathbb{P}}
\newcommand{\R}{\mathbb{R}}
\newcommand{\Z}{\mathbb{Z}}
\newcommand{\N}{\mathbb{N}}
\newcommand{\C}{\mathbb{C}}
\newcommand{\ii}{\mathrm{i}}
\newcommand{\e}{\mathrm{e}}
\newcommand{\dd}{\mathrm{d}}
\newcommand{\one}{\mathbf{1}}
\newcommand{\Tr}{\operatorname{Tr}}
\newcommand{\tr}{\operatorname{tr}}
\newcommand{\Var}{\operatorname{Var}}
\newcommand{\supp}{\operatorname{supp}}
\newcommand{\diag}{\operatorname{diag}}
\newcommand{\dist}{\operatorname{dist}}
\newcommand{\abs}[1]{\left|#1\right|}
\newcommand{\norm}[1]{\left\|#1\right\|}
\newcommand{\ket}[1]{|#1\rangle}
\newcommand{\bra}[1]{\langle #1|}
\newcommand{\braket}[2]{\langle #1|#2\rangle}
\newcommand{\avg}[1]{\left\langle #1\right\rangle}
\newcommand{\dGamma}{\mathrm{d}\Gamma}
\newcommand{\cH}{\mathcal{H}}
\newcommand{\cK}{\mathcal{K}}
\newcommand{\cF}{\mathcal{F}}
\newcommand{\cU}{\mathcal{U}}
\newcommand{\hq}{h_{q}}          % one-body dual Hamiltonian in sector q
\newcommand{\Qc}{Q}              % eigenfunction correlator
\newcommand{\Sfull}{S_{\mathrm{full}}}
\newcommand{\SG}{S_{\mathrm{G}}}
\newcommand{\Smix}{S_{\mathrm{mix}}}
\newcommand{\SN}{S_{N}}

% Section-local macros of the paper, collected by make_statements.py.
\providecommand{\sxb}[1]{\sigma^{x}_{b_{#1}}}
\providecommand{\szb}[1]{\sigma^{z}_{b_{#1}}}
\providecommand{\taux}[1]{\tau^{x}_{#1}}
\providecommand{\tauz}[1]{\tau^{z}_{#1}}
\providecommand{\cdr}{\tilde c}
\providecommand{\LinkS}{\mathcal{S}}
\providecommand{\SiteT}{\mathcal{T}}
\providecommand{\Ilam}{I_\lambda}
\providecommand{\gstar}{\gamma_*}
\providecommand{\Bad}{\mathcal{B}}
\providecommand{\Ends}{\mathcal{E}}
\providecommand{\Wr}{\mathcal{W}}
\providecommand{\sgnob}{\operatorname{sgn}}
\providecommand{\eone}{e_1}
\providecommand{\bent}{\mathsf{h}}
\providecommand{\Hsh}{\mathrm{H}}
\providecommand{\Sfm}{S_{f}}
\providecommand{\Ibar}{\bar I}
\providecommand{\rhoqf}[1]{\rho[#1]}
\providecommand{\Alg}{\mathfrak{A}}
\providecommand{\BAlg}{\mathfrak{B}}
\providecommand{\GAlg}{\mathfrak{G}}
\providecommand{\SigF}{\Sigma_1(F)}
\providecommand{\SigpF}{\Sigma_1'(F)}
\providecommand{\ThF}{\Theta(F)}
\providecommand{\FZ}{F_{\Z}}
\providecommand{\Zone}{\one_N}
\providecommand{\Edim}{\E_{\mathrm{dim}}}
\providecommand{\mudim}{\mu_{\mathrm{dim}}}
\providecommand{\Psidim}{\Psi^{\mathrm{dim}}}
\providecommand{\chidim}{\chi^{\mathrm{dim}}}
\providecommand{\Omdim}{\Omega^{\mathrm{dim}}}
\providecommand{\ncdw}{n^{\mathrm{cdw}}}
\providecommand{\Tdim}{\mathcal{T}}
\providecommand{\Mom}{\mathcal{M}}
\providecommand{\Rot}{\mathrm{Rot}}
\providecommand{\gdim}{\mathsf{g}}
\providecommand{\Cov}{\operatorname{Cov}}
\providecommand{\TLat}{\mathcal{L}}
\providecommand{\TCyc}{\mathcal{Z}}
\providecommand{\TFlux}{\Phi}
\providecommand{\TGauge}{\Theta}
\providecommand{\Tpi}{\pi}
\providecommand{\Tclus}{\mathcal{C}}
\providecommand{\Ush}{\mathrm{H}}
\providecommand{\Uodd}{o}
\providecommand{\UQ}{\mathsf{Q}}
\providecommand{\Femp}{F_{\mathrm{emp}}}
\providecommand{\kaphat}{\hat\kappa}
\providecommand{\psicdw}{\psi_{\mathrm{CDW}}}


% Statements: heading "Theorem 4.7 (Title)." with the paper's number, read from ../main.aux.
\newcommand{\stmtname}{}
\newtheoremstyle{stmtplain}{12pt}{3pt}{\itshape}{}{\bfseries}{.}{.5em}{\thmname{#1}\thmnote{\ \normalfont(#3)}}
\newtheoremstyle{stmtdef}{12pt}{3pt}{\normalfont}{}{\bfseries}{.}{.5em}{\thmname{#1}\thmnote{\ \normalfont(#3)}}
\theoremstyle{stmtplain}
\newtheorem*{stmtP}{\stmtname}
\theoremstyle{stmtdef}
\newtheorem*{stmtD}{\stmtname}
\newcommand{\leanline}[1]{\par\nopagebreak\vspace{1pt}%
  {\raggedright\footnotesize\textsf{Lean:}\ #1\par}\medskip}

\renewcommand{\contentsname}{Indice}

\title{Rigorous disorder-free localization in a $\mathbb{Z}_2$ lattice gauge theory\\[8pt]
\large Sintesi del paper e raccolta degli enunciati}
\author{}
\date{23 novembre 2026}

\begin{document}
\maketitle

\noindent Questo documento riassume il paper (Sezione~\ref{sec:sintesi}) e ne raccoglie tutti gli enunciati
(Sezione~\ref{sec:enunciati}), con la numerazione del paper e, per ciascuno, i teoremi Lean che lo formalizzano.

\tableofcontents

% ------------------------------------------------------------------------------------------------
\begingroup
\hyphenpenalty=10000 \exhyphenpenalty=10000 \emergencystretch=4em
\section{Sintesi}\label{sec:sintesi}

\paragraph{Il problema.}
Nei sistemi quantistici a molti corpi isolati ci si aspetta termalizzazione. La localizzazione a molti corpi (MBL)
propone che un forte disordine congelato la impedisca anche in presenza di interazioni, ma il suo status è
dibattuto e i risultati rigorosi sono pochi, tutti per Hamiltoniane con disordine. Smith, Knolle, Kovrizhin e
Moessner (SKKM) hanno trovato numericamente un meccanismo in cui l'invarianza di gauge prende il posto del
disordine. In una teoria di gauge $\Z_2$ unidimensionale con fermioni liberi, Hamiltoniana e stato iniziale sono
invarianti per traslazione, eppure la memoria di un'onda di densità iniziale persiste e le correlazioni non si
propagano: è la \emph{disorder-free localization} (DFL). Finora l'evidenza era solo numerica. Il paper ne dà una
trattazione rigorosa.

\paragraph{Modello e dualità (Sezione~\ref{sec:model}).}
Fermioni senza spin saltano su una catena di $N$ siti, accoppiati minimamente a spin di Ising sui link, con un
termine elettrico che accoppia link vicini (Definizione~\ref{def:model}). L'unica modifica rispetto a SKKM è un
\emph{link di bordo}, un flusso $\Z_2$ statico all'estremo destro della catena. Con esso una dualità di tipo
Kramers--Wannier è un'equivalenza unitaria esatta sull'intero spazio di Hilbert (Teorema~\ref{thm:duality}):
\[
  U^\dagger HU=\bigoplus_q\Bigl(\dGamma(\hq)-h\sum_jq_j\Bigr),\qquad \hq=-J\Delta_N+2h\diag(q),
\]
dove $q\in\{\pm1\}^N$ sono le cariche di Gauss, conservate. Per lo stato di gauge polarizzato lungo $z$ le cariche
sono esattamente i.i.d.: ogni valore di aspettazione di osservabili gauge-invarianti, a qualunque numero di tempi,
è una media su catene di Anderson con potenziale di Bernoulli $\pm2h$ (Corollario~\ref{cor:sector}). La dualità è
locale attraverso un taglio a meno di una sola porta controllata, e cambia le entropie di entanglement al più di
$\ln2$ (Lemma~\ref{lem:cutlocal}).

\paragraph{L'input a un corpo (Sezione~\ref{sec:onebody}).}
Tutti i risultati sulla teoria di gauge discendono da un'unica ipotesi sulla catena di Bernoulli,
l'Ipotesi~\ref{hyp:EC}: il correlatore delle autofunzioni, mediato sulle cariche, decade in modo sommabile,
uniformemente nel volume. Il Teorema~\ref{thm:EC} la dimostra con decadimento esponenziale,
$\E\,Q^\Lambda(x,y)\le Ce^{-\eta|x-y|}$, per ogni intervallo finito $\Lambda$ e ogni coppia $x,y\in\Lambda$, anche
al bordo, e anche con l'estremo superiore su $\Lambda$ dentro l'attesa. Le costanti dipendono solo da $|\lambda|$,
con $\lambda=2h/J$, e non sono esplicite. La dimostrazione adatta l'argomento di risonanza a due punti di Macera e
l'eliminazione delle risonanze di Macera e Sodin; i box troncati dal bordo sono trattati con una stima
deterministica delle funzioni di Green (Lemma~\ref{lem:green}). L'input dalla letteratura sono le grandi deviazioni
uniformi per prodotti di matrici di trasferimento di Bernoulli (Proposizione~\ref{prop:LDT}, teoria di
Furstenberg), le cui ipotesi sono verificate nella Proposizione~\ref{prop:nocritical}.

\paragraph{Localizzazione senza disordine della teoria di gauge (Sezione~\ref{sec:main}).}
Per ogni $h\neq0$, uniformemente nel volume $N$ e nel tempo:
\begin{itemize}[leftmargin=1.5em]
\item i fermioni vestiti sono localizzati e la carica trasportata attraverso un taglio qualsiasi è limitata
  (Teorema~\ref{thm:transport});
\item vale un bound di Lieb--Robinson a velocità zero in media sui settori, per ogni stato che dà lo stesso peso a
  tutti i settori; non esiste una versione in norma operatoriale, perché il settore pulito $q\equiv1$ è balistico
  (Proposizione~\ref{prop:LR});
\item l'imbalance a tempo infinito di un'onda di densità iniziale è strettamente positivo
  (Teorema~\ref{thm:imbalance}), anche sulla catena aperta di SKKM senza link di bordo
  (Corollario~\ref{cor:imbalanceopen});
\item l'entanglement condizionato alle cariche, $\SG$, e, in media sulle cariche, l'entanglement completo materia
  più gauge degli stati iniziali a carica definita soddisfano leggi d'area (Teorema~\ref{thm:arealaw}); la
  Proposizione~\ref{prop:entropies} stabilisce le relazioni fra le entropie: l'entropia di numero $\SN$ è un
  minorante di $\Sfull$, di $\Smix$ e dell'entropia della materia nuda, ma non di $\SG$;
\item l'entropia di numero è limitata, con costante esplicita (Teorema~\ref{thm:numberentropy});
\item i limiti termodinamici esistono su catene infinite in entrambe le direzioni, con le stesse stime per le
  osservabili di bulk, imbalance escluso (Teorema~\ref{thm:infvol}).
\end{itemize}
Le costanti sono esplicite in termini del profilo di decadimento del Teorema~\ref{thm:EC}. Non è controllato
l'entanglement $\Sfull$ dello stato sovrapposto, che dipende dalle coerenze fra settori. Il paper congettura che
cresca logaritmicamente (Congettura~\ref{conj:logent}) e ne isola l'ingrediente rigoroso, una stima di sensibilità
alle cariche lontane (Lemma~\ref{lem:duhamel}); congettura inoltre una legge di volume per le entropie miste mediate
nel tempo (Congettura~\ref{conj:volmix}).

\paragraph{La localizzazione dipende dallo stato iniziale (Sezione~\ref{sec:dimer}).}
Cambiando solo lo stato dei link pari si ottiene uno stato prodotto la cui legge dei settori è la legge di dimero
casuale (Definizione~\ref{def:dimerstate}, Proposizione~\ref{prop:dimerlaw}). Per $0<2|h|<J$ il modello di dimero
ha energie critiche (Lemma~\ref{lem:dimercritical}). Un fermione vestito si propaga allora almeno in modo
diffusivo, in ogni settore e in ogni volume, fino a tempi dell'ordine della distanza dal bordo
(Proposizione~\ref{prop:dimerlower}), e in modo superdiffusivo nel limite termodinamico, con spostamento quadratico
medio almeno dell'ordine di $T^{3/2}$ in media (Proposizione~\ref{prop:dimersuper}, Teorema~\ref{thm:dimer}). Per
ogni $h\neq0$ la media a tempo infinito dell'imbalance tende a zero per $N\to\infty$
(Proposizione~\ref{prop:dimerimbalance}).

\paragraph{Due dimensioni (Sezione~\ref{sec:2d}).}
La mappa è esatta su ogni grafo finito, con l'insieme completo delle etichette di superselezione, compresi i loop di
Wilson non contrattili di un toro (Lemma~\ref{lem:2dgauge}, Teorema~\ref{thm:2dmap}). Gli stati prodotto portano al
modello di Anderson bidimensionale con potenziale di Bernoulli, oppure a un flusso $\pi$ casuale binario
(Proposizione~\ref{prop:2dlaws}). Ad accoppiamento forte le cariche restano confinate in cluster finiti fino a tempi
dell'ordine di $|h|/J^2$ (Proposizione~\ref{prop:cluster}). La localizzazione in due dimensioni non è dimostrata:
manca l'input a un corpo.

\paragraph{Teorie di gauge $U(1)$ (Sezione~\ref{sec:u1}).}
Con link quantistici di spin $S$, la legge di Gauss limita l'entropia di numero di ogni intervallo con
$\ln(4S+1)$, in ogni settore di carica di fondo, per ogni stato e ogni dinamica gauge-invariante
(Proposizione~\ref{prop:gausscap}). Con link di Wilson, la conservazione dell'energia ne limita la media sui tagli
(Proposizione~\ref{prop:wilson}). Per gli stati iniziali di Brenes et al.\ i campi di fondo sono i.i.d.\ e agiscono
sulle cariche come un potenziale a cammino aleatorio (Proposizione~\ref{prop:brenes}); nella numerica esatta la loro
memoria a tempi lunghi supera quella dovuta all'energia solo ad accoppiamento forte (Osservazione~\ref{rem:brenes}).
Nella teoria di gauge $\Z_2$ standard, con termine elettrico a un solo link, una
carica si muove in un potenziale di cammino aleatorio e due cariche interagiscono.

\paragraph{Numerica e problemi aperti (Sezioni~\ref{sec:numerics} e~\ref{sec:discussion}).}
Simulazioni esatte illustrano i risultati e confrontano le costanti rigorose con i dati. Per $N=10$ l'entanglement
completo cresce per diverse decadi di tempo, in accordo con la Congettura~\ref{conj:logent}. Il principale problema
aperto è la materia interagente: con un'interazione di densità la media sui settori resta esatta, e la
localizzazione senza disordine diventa MBL di una catena XXZ in un campo casuale binario, fuori dalla portata dei
metodi rigorosi attuali.

\paragraph{Verifica formale in Lean (Appendice~\ref{app:lean}).}
I risultati del paper, in volume finito e infinito, sono formalizzati in Lean~4 con Mathlib: 96 moduli (circa
51\,800 righe) e 5752 teoremi, senza assiomi oltre ai tre standard di Lean. Sono formalizzate anche tutte le
osservazioni con contenuto matematico, comprese la stima (EC) sull'anello, la localizzazione del fermione vestito
ad accoppiamento forte e la riduzione esatta del caso interagente. Restano fuori solo i risultati di letteratura usati come ipotesi (grandi deviazioni di
Furstenberg, Ge--Zhao, la stima di percolazione di Li--Liang, la localizzazione dinamica di De Bièvre--Germinet) e
la verifica delle loro ipotesi nel formalismo originale.
\endgroup

% ------------------------------------------------------------------------------------------------
\clearpage
\section{Enunciati}\label{sec:enunciati}

\noindent Seguono tutti gli enunciati del paper (definizioni, ipotesi, lemmi, proposizioni, teoremi, corollari e
congetture: 58 in tutto), nell'ordine e con la numerazione del paper. Il testo è quello del paper, in inglese; i
rimandi a equazioni, osservazioni, figure e bibliografia si riferiscono al paper. Sotto ogni enunciato sono elencati
i teoremi Lean che lo formalizzano, con i file della libreria (\texttt{lean/DFL/}), secondo la tabella di
corrispondenza del README della formalizzazione. Le dimostrazioni non sono riportate.

% Generated by make_statements.py -- do not edit by hand.
\subsection*{\S\,\ref{sec:model}\quad The model and its exact free-fermion duality}
\addcontentsline{toc}{subsection}{\S\,\ref{sec:model}\quad The model and its exact free-fermion duality}
\renewcommand{\stmtname}{\Cref{def:model}}
\begin{stmtD}[The model]
Let $J>0$ and $h\in\R\setminus\{0\}$.  The Hamiltonian on $\cH_N$ is
\begin{equation}\tag{\ref{eq:H}}
  H=-J\sum_{j=1}^{N-1}\szb{j}\bigl(f_j^\dagger f_{j+1}+f_{j+1}^\dagger f_j\bigr)
    -h\sum_{j=1}^{N}\sxb{j-1}\,\sxb{j}\qquad(\sxb{0}:=\one),
\end{equation}
and the Gauss operators are
\begin{equation}\tag{\ref{eq:gauss}}
  G_j:=\sxb{j-1}\,(-1)^{n_j}\,\sxb{j},\qquad j=1,\dots,N .
\end{equation}
An operator $O$ on $\cH_N$ is \emph{gauge invariant} if $[O,G_j]=0$ for all $j$.  For
$q\in\{\pm1\}^N$ let $P_q:=\prod_{j=1}^N\frac12(\one+q_jG_j)$.
\end{stmtD}
\leanline{\texttt{H}, \texttt{gauss}, \texttt{ann}, \texttt{cre}, \texttt{dGamma}, \texttt{hq}, \texttt{Hsec}, \texttt{Hdual} \,(\texttt{Basic.\allowbreak{}lean})}
\renewcommand{\stmtname}{\Cref{lem:gauss}}
\begin{stmtP}[Gauss law]
\leavevmode
\begin{enumerate}[label=(\roman*)]
\item The $G_j$ are self-adjoint unitaries, commute with each other, with $N_f$ and with $H$.
\item For every $j\in\{1,\dots,N\}$ one has the operator identity (one-dimensional $\Z_2$ Gauss law)
\begin{equation}\tag{\ref{eq:gausslaw1d}}
  \sxb{j}=(-1)^{N_{\le j}}\prod_{k\le j}G_k .
\end{equation}
In particular $\prod_{j=1}^NG_j=(-1)^{N_f}\sxb{N}$, and the boundary flux $\sxb{N}$ commutes with $H$.
\item For every $q\in\{\pm1\}^N$ the projection $P_q$ has rank $2^N$; thus $\cH_N=\bigoplus_q
P_q\cH_N$ with $2^N$ non-trivial superselection sectors, each invariant under $H$.
\end{enumerate}
\end{stmtP}
\leanline{\texttt{gauss\_\allowbreak{}isSymm}, \texttt{gauss\_\allowbreak{}mul\_\allowbreak{}self}, \texttt{gauss\_\allowbreak{}comm}, \texttt{gauss\_\allowbreak{}comm\_\allowbreak{}numF}, \texttt{gauss\_\allowbreak{}comm\_\allowbreak{}H}, \texttt{gauss\_\allowbreak{}law\_\allowbreak{}1d}, \texttt{gauss\_\allowbreak{}prod}, \texttt{sxL\_\allowbreak{}last\_\allowbreak{}comm\_\allowbreak{}H}, \texttt{gauss\_\allowbreak{}mul\_\allowbreak{}gaussProj}, \texttt{gaussProj\_\allowbreak{}mul}, \texttt{sum\_\allowbreak{}gaussProj}, \texttt{rank\_\allowbreak{}gaussProj}, \texttt{gaussProj\_\allowbreak{}comm\_\allowbreak{}H} \,(\texttt{GaussLaw.\allowbreak{}lean})}
\renewcommand{\stmtname}{\Cref{thm:duality}}
\begin{stmtP}[Exact duality]
$U$ is unitary, preserves all occupation numbers (in particular it is even), and:
\begin{enumerate}[label=(\alph*)]
\item \emph{Gauss operators:} $U^\dagger G_jU=\taux{j}$ for all $j$; hence
$U^\dagger P_qU=\one\otimes\ket{q}\bra{q}$.
\item \emph{Hamiltonian:}
\begin{equation}\tag{\ref{eq:dualH}}
  U^\dagger HU=-J\sum_{j=1}^{N-1}(c_j^\dagger c_{j+1}+c_{j+1}^\dagger c_j)-h\sum_{j=1}^N\taux{j}(-1)^{n_j}
  =\sum_{q\in\{\pm1\}^N}\Bigl(\dGamma(\hq)-h\sum_{j=1}^Nq_j\Bigr)\otimes\ket{q}\bra{q}.
\end{equation}
Consequently, with $E_1(q)<\dots<E_N(q)$ the eigenvalues of $\hq$ (simple, since $J\neq0$ and $\hq$ is a
Jacobi matrix with non-zero off-diagonal entries), the spectrum of $H$ (with
multiplicities) is the union over $q$ of the multisets
$\{\sum_\alpha m_\alpha E_\alpha(q)-h\sum_jq_j:\ m\in\{0,1\}^N\}$.
\item \emph{Dictionary:} for all $j$ (with $\tauz{N+1}=\one$),
\begin{gather*}
  U^\dagger n_jU=n_j,\qquad U^\dagger(-1)^{N_f}U=(-1)^{N_c},\qquad U^\dagger f_jU=c_j\tauz{j},\\
  U^\dagger\szb{j}U=\tauz{j}\tauz{j+1},\qquad
  U^\dagger\Bigl(\prod_{k=j}^N\szb{k}\Bigr)U=\tauz{j},\\
  U^\dagger\sxb{j}U=\prod_{k\le j}\taux{k}(-1)^{n_k},\qquad
  U^\dagger\sxb{j-1}\sxb{j}U=\taux{j}(-1)^{n_j},\\
  U^\dagger\bigl(\szb{j}f_j^\dagger f_{j+1}\bigr)U=c_j^\dagger c_{j+1}\quad(j<N).
\end{gather*}
\item \emph{Dressed fermions:} the operators
\begin{equation}\tag{\ref{eq:dressed}}
  \cdr_j:=Uc_jU^\dagger=f_j\prod_{k=j}^{N}\szb{k}
\end{equation}
satisfy the CAR, commute with every $G_k$, are odd, and $\cdr_j^\dagger\cdr_j=n_j$.  For $x<y$,
$\cdr_x^\dagger\cdr_y=f_x^\dagger\bigl(\prod_{x\le k<y}\szb{k}\bigr)f_y$, the gauge-invariant
hopping with a Wilson line between $x$ and $y$.  In terms of the original operators,
\begin{equation}\tag{\ref{eq:Hdressed}}
  H=-J\sum_{j=1}^{N-1}(\cdr_j^\dagger\cdr_{j+1}+\cdr_{j+1}^\dagger\cdr_j)-h\sum_{j=1}^NG_j(-1)^{n_j}
   =\sum_q P_q\Bigl(\sum_{x,y}(\hq)_{xy}\cdr_x^\dagger\cdr_y-h\sum_jq_j\Bigr).
\end{equation}
\item \emph{Dynamics:}
$U^\dagger\e^{-\ii tH}U=\sum_q\e^{\ii ht\sum_jq_j}\,\Gamma(\e^{-\ii t\hq})\otimes\ket{q}\bra{q}$, and
\begin{equation}\tag{\ref{eq:heisenberg}}
  \e^{\ii tH}\cdr_x\e^{-\ii tH}=\sum_qP_q\sum_{y=1}^N(\e^{-\ii t\hq})_{xy}\,\cdr_y\qquad(t\in\R).
\end{equation}
\end{enumerate}
\end{stmtP}
\leanline{\texttt{duality}, \texttt{gauss\_\allowbreak{}dual}, \texttt{dGamma\_\allowbreak{}hq}, \texttt{hop\_\allowbreak{}apply}, \texttt{numOp\_\allowbreak{}dual} \,(\texttt{Duality.\allowbreak{}lean});\\ \texttt{ann\_\allowbreak{}dual}, \texttt{sz\_\allowbreak{}dual}, \texttt{szTail\_\allowbreak{}dual}, \texttt{sx\_\allowbreak{}dual}, \texttt{elec\_\allowbreak{}dual}, \texttt{hopTerm\_\allowbreak{}dual}, \texttt{parity\_\allowbreak{}dual}, \texttt{dressed\_\allowbreak{}eq}, \texttt{dressed\_\allowbreak{}anticomm}, \texttt{dressed\_\allowbreak{}CAR}, \texttt{dressed\_\allowbreak{}comm\_\allowbreak{}gauss}, \texttt{dressed\_\allowbreak{}odd}, \texttt{dressed\_\allowbreak{}number}, \texttt{dressed\_\allowbreak{}hop}, \texttt{H\_\allowbreak{}dressed}, \texttt{evo\_\allowbreak{}dual}, \texttt{heisenberg\_\allowbreak{}dressed} \,(\texttt{DualityFull.\allowbreak{}lean});\\ \texttt{charpoly\_\allowbreak{}dGamma}, \texttt{charpoly\_\allowbreak{}H} \,(\texttt{DualitySpectrum.\allowbreak{}lean})}
\renewcommand{\stmtname}{\Cref{prop:variants}}
\begin{stmtP}[Open and ring variants]
\leavevmode
\begin{enumerate}[label=(\roman*)]
\item \emph{Open chain.}  Let $H^{\rm open}$ and $G^{\rm open}_j$ be given by \eqref{eq:H},
\eqref{eq:gauss} on $\cF_N\otimes(\C^2)^{\otimes(N-1)}$ with $\sxb{0}=\sxb{N}:=\one$, and let
$\iota\chi:=\chi\otimes\ket{+}_{b_N}$.  Then $H(\one\otimes\iota)=(\one\otimes\iota)H^{\rm open}$ and
$G_j(\one\otimes\iota)=(\one\otimes\iota)G_j^{\rm open}$, so $H^{\rm open}$ is unitarily equivalent to
$H$ restricted to the invariant subspace $\{\sxb{N}=+1\}$, and
\begin{equation}\tag{\ref{eq:openparity}}
  U^\dagger\tfrac12(\one+\sxb{N})U=\sum_q\Pi_{\pi(q)}\otimes\ket{q}\bra{q},\qquad\pi(q):=\prod_{j}q_j,
\end{equation}
where $\Pi_{\pm}$ is the projection onto $\{(-1)^{N_c}=\pm1\}$ in $\cF^c_N$.  Thus
$H^{\rm open}\cong\bigoplus_q\bigl(\dGamma(\hq)-h\sum_jq_j\bigr)\big|_{\Pi_{\pi(q)}\cF^c_N}$ (the
parity constraint $\prod_jq_j=(-1)^{N_f}$).  For every unit vector $\psi\in\cF_N$, the state
$\psi\otimes\ket{\uparrow\cdots\uparrow}_z\in\cF_N\otimes(\C^2)^{\otimes(N-1)}$ satisfies
\begin{equation}\tag{\ref{eq:openinit}}
  U^\dagger(\one\otimes\iota)\bigl(\psi\otimes\ket{\uparrow\cdots\uparrow}_z\bigr)
  =2^{-(N-1)/2}\sum_q\Pi_{\pi(q)}\psi\otimes\ket{q}.
\end{equation}
So its sector law $\mu(q)=2^{-(N-1)}\norm{\Pi_{\pi(q)}\psi}^2$ has density
$2\norm{\Pi_{\pi(q)}\psi}^2\le2$ with respect to the i.i.d.\ Bernoulli$(1/2)$ law, and the sector $q$
carries the matter state $\Pi_{\pi(q)}\psi$ (normalized).  If $\psi$ has definite parity $\pi$, the sector
law is uniform on $\{q:\prod_jq_j=\pi\}$, with density $2\cdot\one[\prod_jq_j=\pi]$, and every sector in
its support carries the same matter state $\psi$.
\item \emph{Ring} ($N\ge3$).  Let $H^{\rm ring}$ be \eqref{eq:H} on $\cH_N$ with the hopping sum
extended to $j=N$ ($f_{N+1}:=f_1$) and $\sxb{0}:=\sxb{N}$, let $G^{\rm ring}_j$ be \eqref{eq:gauss}
with $\sxb{0}:=\sxb{N}$, and $W:=\prod_{j=1}^N\szb{j}$.  Then, with the same $U$,
\begin{gather*}
  U^\dagger G^{\rm ring}_1U=(-1)^{N_c}\prod_{k\ge2}\taux{k},\qquad U^\dagger G^{\rm ring}_jU=\taux{j}\ (j\ge2),
  \qquad U^\dagger WU=\tauz{1},\\
  \begin{aligned}
  U^\dagger H^{\rm ring}U={}&-J\sum_{j=1}^{N-1}(c_j^\dagger c_{j+1}+{\rm h.c.})-J\tauz{1}(c_N^\dagger c_1+c_1^\dagger c_N)\\
  &-h(-1)^{N_c}\Bigl(\prod_{k\ge2}\taux{k}\Bigr)(-1)^{n_1}-h\sum_{j=2}^N\taux{j}(-1)^{n_j}.
  \end{aligned}
\end{gather*}
Consequently $W$ and the $G^{\rm ring}_j$ commute with $H^{\rm ring}$, $\prod_jG_j^{\rm ring}=(-1)^{N_f}$, and on
the joint eigenspace $\{G^{\rm ring}_j=q_j\ \forall j,\ W=w\}$ (non-zero for every $q$ and $w$) the
Hamiltonian $H^{\rm ring}$ is unitarily equivalent to $\dGamma(h_{q,w})-h\sum_jq_j$ restricted to
$\Pi_{\pi(q)}\cF^c_N$, where
\begin{equation}\tag{\ref{eq:ringtwist}}
  h_{q,w}:=\hq-Jw\bigl(\ket{1}\bra{N}+\ket{N}\bra{1}\bigr).
\end{equation}
The twist on the bond $(N,1)$ is exactly $w$, the eigenvalue of the Wilson loop, with no
parity-dependent factor.  The state $\psi\otimes\ket{\uparrow\cdots\uparrow}_z$ with $\psi$ of parity
$\pi$ has $W=+1$ (periodic $h_{q,1}$) and sector law uniform on $\{q:\prod_jq_j=\pi\}$.
\end{enumerate}
\end{stmtP}
\leanline{\texttt{H\_\allowbreak{}boundary\_\allowbreak{}link}, \texttt{linkOf\_\allowbreak{}last}, \texttt{sector\_\allowbreak{}average\_\allowbreak{}open} \,(\texttt{Variants.\allowbreak{}lean});\\ \texttt{ring\_\allowbreak{}duality}, \texttt{gaussRing\_\allowbreak{}dual\_\allowbreak{}zero}, \texttt{gaussRing\_\allowbreak{}dual}, \texttt{wilson\_\allowbreak{}dual}, \texttt{Hring\_\allowbreak{}comm\_\allowbreak{}gaussRing}, \texttt{Hring\_\allowbreak{}comm\_\allowbreak{}wilson}, \texttt{gaussRing\_\allowbreak{}prod}, \texttt{ringVec\_\allowbreak{}orthonormal}, \texttt{card\_\allowbreak{}ringLabels}, \texttt{Hring\_\allowbreak{}ringVec}, \texttt{sector\_\allowbreak{}average\_\allowbreak{}ring} \,(\texttt{Ring.\allowbreak{}lean})}
\renewcommand{\stmtname}{\Cref{lem:cutlocal}}
\begin{stmtP}[Cut-locality]
\leavevmode
\begin{enumerate}[label=(\alph*)]
\item $U=C_\ell\,(U_{\le\ell}\otimes U_{>\ell})$.
\item For every unit vector $\Phi\in\cK_N$: $\ \abs{S_A(U\Phi)-S_{A'}(\Phi)}\le\ln2$.
\item If $\Phi$ is an eigenvector of $(-1)^{N_c}$ and of $\prod_{k\le\ell}\taux{k}$ (which is
$U^\dagger\prod_{k\le\ell}G_kU$, the total $\Z_2$ charge in $A$), then $S_A(U\Phi)=S_{A'}(\Phi)$.  In
particular, for a definite-charge state $\Phi=\phi\otimes\ket{q}$ with $\phi\in\cF_N^c$ of definite
parity, $S_A(U\Phi)=S(\rho^\phi_{\le\ell})$, the entanglement entropy of $\phi$ for the modes
$1,\dots,\ell$.
\item The constant $\ln2$ in (b) is attained in both directions (already for $N=2$, $\ell=1$).
\end{enumerate}
\end{stmtP}
\leanline{\texttt{vonNeumann\_\allowbreak{}rdmLGT\_\allowbreak{}xprod}, \texttt{lgt\_\allowbreak{}areaLaw\_\allowbreak{}full} \,(\texttt{CutLocal.\allowbreak{}lean});\\ \texttt{EntropiesAux.\allowbreak{}dualEquiv\_\allowbreak{}symm\_\allowbreak{}splitLGT}, \texttt{cutlocal\_\allowbreak{}defect}, \texttt{exProd\_\allowbreak{}entropy}, \texttt{exProd\_\allowbreak{}U\_\allowbreak{}entropy}, \texttt{exProd\_\allowbreak{}Udag\_\allowbreak{}entropy} \,(\texttt{Entropies.\allowbreak{}lean}, \texttt{CutLocalExamples.\allowbreak{}lean})}
\renewcommand{\stmtname}{\Cref{lem:pinching}}
\begin{stmtP}[Two-block pinching]
Let $P^+$, $P^-$ be orthogonal projections on a finite-dimensional space with $P^++P^-=\one$, and let
$\mathcal{P}(\sigma):=P^+\sigma P^++P^-\sigma P^-$.  Then
$S(\sigma)\le S(\mathcal{P}(\sigma))\le S(\sigma)+\ln2$ for every density matrix $\sigma$.
\end{stmtP}
\leanline{\texttt{vonNeumann\_\allowbreak{}le\_\allowbreak{}pinching}, \texttt{vonNeumann\_\allowbreak{}pinching\_\allowbreak{}le}, \texttt{vonNeumann\_\allowbreak{}concave}, \texttt{vonNeumann\_\allowbreak{}mix\_\allowbreak{}le}, \texttt{vonNeumann\_\allowbreak{}le\_\allowbreak{}diag}, \texttt{vonNeumann\_\allowbreak{}le\_\allowbreak{}log\_\allowbreak{}card}, \texttt{vonNeumann\_\allowbreak{}mul\_\allowbreak{}conjTranspose\_\allowbreak{}comm} \,(\texttt{VNEntropy.\allowbreak{}lean})}
\renewcommand{\stmtname}{\Cref{cor:sector}}
\begin{stmtP}[Exact sector averaging]
Let $\psi\in\cF_N$ be a unit vector, $\Psi_0:=\ket{\uparrow\cdots\uparrow}_z\otimes\psi$ (we do not
distinguish $\LinkS_N\otimes\cF_N$ from $\cF_N\otimes\LinkS_N$), and $\Psi(t):=\e^{-\ii tH}\Psi_0$.
\begin{enumerate}[label=(\alph*)]
\item One has
\[
  U^\dagger\Psi_0=2^{-N/2}\sum_q\psi\otimes\ket{q},\qquad
  U^\dagger\Psi(t)=2^{-N/2}\sum_q\e^{\ii ht\sum_jq_j}\,\psi_q(t)\otimes\ket{q}.
\]
In particular $\norm{P_q\Psi_0}^2=2^{-N}$ for every $q$: the charges are exactly i.i.d.\
Bernoulli$(1/2)$, and every sector carries the same matter state $\psi$.
\item For every gauge-invariant $O$ and $t\in\R$,
\begin{equation}\tag{\ref{eq:sectoravg}}
  \langle\Psi(t),O\,\Psi(t)\rangle=\E_q\,\langle\psi_q(t),O_q\,\psi_q(t)\rangle .
\end{equation}
More generally, for gauge-invariant $O_1,\dots,O_m$ and $t_1,\dots,t_m\in\R$,
\[
  \langle\Psi_0,O_1(t_1)\cdots O_m(t_m)\Psi_0\rangle=\E_q\langle\psi,O_{1,q}(t_1)\cdots O_{m,q}(t_m)\psi\rangle,
\]
where $O(t)=\e^{\ii tH}O\e^{-\ii tH}$.  The constants $-h\sum_jq_j$ drop out.
\item \emph{Sector-diagonal observables.}  The following are gauge invariant, with the stated $O_q$:
\begin{itemize}
\item every function $F(n_1,\dots,n_N)$ of the occupations, with $O_q=F(n_1,\dots,n_N)$
independent of $q$; in particular $n_x$ and $\one[N_X=k]$.  Hence the distribution of $N_X$ in
$\Psi(t)$ is $\Prob_{\Psi(t)}(N_X=k)=\E_q\langle\psi_q(t),\one[N_X=k]\psi_q(t)\rangle$, the
$q$-average of the per-sector distributions;
\item every polynomial in the dressed fermions $\cdr_x,\cdr_x^\dagger$, with $O_q$ the same
polynomial in $c_x,c_x^\dagger$; in particular $f_x^\dagger(\prod_{x\le k<y}\szb{k})f_y\mapsto c_x^\dagger c_y$;
\item the electric field and energy: $\sxb{j}\mapsto(\prod_{k\le j}q_k)(-1)^{N_{\le j}}$ and
$\sxb{j-1}\sxb{j}\mapsto q_j(-1)^{n_j}$, together with their products.
\end{itemize}
For a Fock state $\psi=\ket{n}$ with occupied set $\{s_1<\dots<s_M\}$, write $U_q(t):=\e^{-\ii t\hq}$ (a unitary $N\times N$ matrix; the subscript distinguishes it from the duality unitary $U$).
Then $\langle\psi_q(t),c_x^\dagger c_y\psi_q(t)\rangle=\sum_{a=1}^M\overline{U_q(t)_{xs_a}}\,U_q(t)_{ys_a}$, and
\begin{equation}\tag{\ref{eq:stringdet}}
  \langle\psi_q(t),(-1)^{N_{\le j}}\psi_q(t)\rangle=\det\bigl[(U_q(t)^*R_jU_q(t))_{s_as_b}\bigr]_{a,b=1}^M,
  \qquad R_j:=\diag\bigl((-1)^{\one[x\le j]}\bigr)_{x=1}^N .
\end{equation}
\item \emph{The gauge field $\szb{j}$ is not sector diagonal.}  It anticommutes with $G_j$ and
$G_{j+1}$ (with $G_N$ only, if $j=N$), and
\begin{equation}\tag{\ref{eq:loschmidt}}
  \langle\Psi(t),\szb{j}\Psi(t)\rangle=\E_q\Bigl[\e^{\ii ht(\sum_kq_k-\sum_kq'_k)}\,
  \langle\psi_{q'}(t),\psi_q(t)\rangle\Bigr],\qquad q':=q^{[j]},
\end{equation}
where $q^{[j]}$ is $q$ with $q_j$ and $q_{j+1}$ flipped ($q_N$ only, if $j=N$), so that
$\sum_kq_k-\sum_kq'_k=2(q_j+q_{j+1})$ (resp.\ $2q_N$).  For Fock $\psi$,
$\langle\psi_{q'}(t),\psi_q(t)\rangle=\det[(U_{q'}(t)^*U_q(t))_{s_as_b}]_{a,b=1}^M$.  Thus gauge-field
observables are Loschmidt echoes between sectors that differ by a local flip of the potential, and
the sector-dependent phase in \eqref{eq:loschmidt} cannot be omitted.
\end{enumerate}
\end{stmtP}
\leanline{\texttt{sector\_\allowbreak{}average}, \texttt{sector\_\allowbreak{}average\_\allowbreak{}matter}, \texttt{gaugeInv\_\allowbreak{}submatrix}, \texttt{correlator\_\allowbreak{}sector\_\allowbreak{}average} \,(\texttt{Sector.\allowbreak{}lean}, \texttt{Interacting.\allowbreak{}lean})}
\renewcommand{\stmtname}{\Cref{lem:productlaw}}
\begin{stmtP}[Sector laws of product initial states]
Let $\chi=\bigotimes_{j=1}^N\chi_j\in\LinkS_N$ be a product state of the links.  Write
$\chi(x):=\prod_j\langle x_{b_j}|\chi_j\rangle$ for its $\sigma^x$-amplitudes, and
$\Psi_0=\chi\otimes\psi$.
\begin{enumerate}[label=(\alph*)]
\item For every $\psi=\sum_n\psi(n)\ket{n}$, $U^\dagger\Psi_0=\sum_q\psi^{(q)}\otimes\ket{q}$ with
$\psi^{(q)}:=\sum_n\chi(x(n,q))\psi(n)\ket{n}$.
\item \emph{Fock matter.}  If $\psi=\ket{n}$, then $\psi^{(q)}=\chi(x(n,q))\ket{n}$.  The sector law
$\mu(q):=\norm{P_q\Psi_0}^2=\abs{\chi(x(n,q))}^2$ is the image of the product law
$\bigotimes_j\bigl(\abs{\langle+|\chi_j\rangle}^2\delta_{+1}+\abs{\langle-|\chi_j\rangle}^2\delta_{-1}\bigr)$
of $x$ under the bijection $x\mapsto q$, $q_j=(-1)^{n_j}x_{b_{j-1}}x_{b_j}$, $x_{b_0}=1$.  Every
sector in the support carries the matter state $\ket{n}$, up to a phase, and for every gauge-invariant
$O$, $\langle\Psi(t),O\Psi(t)\rangle=\sum_q\mu(q)\langle n,O_q(t)\,n\rangle$.
\item \emph{$\sigma^z$-product gauge states.}  If $\chi=\ket{s}_z$ with $D:=\{j:s_j=-1\}$, then for every
$\psi$
\[
  U^\dagger\Psi_0=2^{-N/2}\sum_q\varepsilon_D(q)\,(W_D\psi)\otimes\ket{q},\qquad
  \varepsilon_D(q):=\prod_{j\in D}\prod_{k\le j}q_k,\quad W_D:=\prod_{j\in D}(-1)^{N_{\le j}} .
\]
So the charges are i.i.d.\ Bernoulli$(1/2)$, and \cref{cor:sector}(b),(c) hold with $\psi$ replaced
by $W_D\psi$.  For Fock $\psi$, $W_D\psi=\pm\psi$, and the gauge-invariant dynamics is identical to
that of $\ket{\uparrow\cdots\uparrow}_z\otimes\psi$.
\item \emph{$\sigma^x$-product gauge states.}  If $\chi=\ket{\xi}_\sigma$, $\xi\in\{\pm1\}^N$, and
$\psi=\ket{n}$, then $U^\dagger\Psi_0=\ket{n}\otimes\ket{q(n,\xi)}$, a single (definite-charge) sector with
$q_j=(-1)^{n_j}\xi_{b_{j-1}}\xi_{b_j}$.
\item \emph{Mixed product states and the dimer law.}  Let $N=2M$ be even.  Let the odd links
$b_{2k-1}$ be in $\sigma^z$-eigenstates and the even links $b_{2k}$ in the $\sigma^x$-eigenstates
$\ket{\xi_{2k}}$, $\xi_{2k}\in\{\pm1\}$, and let $\psi=\ket{n}$.  Then $(q_{2k-1})_{k=1}^M$ are i.i.d.\ uniform
under $\mu$, and $q_{2k}=d_k\,q_{2k-1}$ with the deterministic signs
$d_k:=(-1)^{n_{2k-1}+n_{2k}}\xi_{2k-2}\xi_{2k}$, $\xi_0:=1$.  In particular, for
\[
  \chi_{2k-1}=\ket{\uparrow}_z,\qquad \chi_{2k}=\ket{(-1)^k}\ \ (\sigma^x\text{-eigenvalue }(-1)^k),\qquad
  n_{2k-1}=0,\quad n_{2k}=1,
\]
for $k=1,\dots,M$,
one has, with all amplitudes positive,
\begin{equation}\tag{\ref{eq:dimerstate}}
  U^\dagger\Psi_0=2^{-M/2}\sum_{\eta\in\{\pm1\}^M}\ket{n}\otimes\ket{\eta_1,\eta_1,\eta_2,\eta_2,\dots,\eta_M,\eta_M}.
\end{equation}
The sector law is the \emph{random dimer law}: i.i.d.\ equal pairs $(q_{2k-1},q_{2k})=(\eta_k,\eta_k)$,
$\eta_k$ i.i.d.\ uniform.  The same CDW matter state sits in every sector.  (\Cref{prop:dimerlaw} extends this
to general Fock matter and identifies the states with $d_k\equiv1$.)
\end{enumerate}
\end{stmtP}
\leanline{\texttt{sector\_\allowbreak{}sum\_\allowbreak{}prod}, \texttt{sector\_\allowbreak{}average\_\allowbreak{}prod\_\allowbreak{}fock}, \texttt{sector\_\allowbreak{}average\_\allowbreak{}zprod}, \texttt{expect\_\allowbreak{}xprod\_\allowbreak{}fock} \,(\texttt{ProductLaw.\allowbreak{}lean})}
\subsection*{\S\,\ref{sec:onebody}\quad One-body input: localization for the Bernoulli chain}
\addcontentsline{toc}{subsection}{\S\,\ref{sec:onebody}\quad One-body input: localization for the Bernoulli chain}
\renewcommand{\stmtname}{\Cref{lem:normalization}}
\begin{stmtP}[Normalization]
Let $\lambda=2h/J$ and let $\Gamma:=\diag((-1)^x)_{x=1}^N$ (as in \cref{thm:imbalance}; not to be confused
with the second quantization $\Gamma(\cdot)$ of \cref{sec:model}).  Then $\Gamma\hq\Gamma^{-1}=J\,A_{[1,N]}$.  Consequently
$Q^{[1,N]}_q(x,y)=Q^{[1,N]}(x,y)$ for all $x,y$ and every $q$.  Moreover, the joint law of
$(Q^\Lambda(x,y))_{x,y\in\Lambda}$ is the same for $\lambda$ and $-\lambda$.
\end{stmtP}
\leanline{\texttt{stag\_\allowbreak{}conj\_\allowbreak{}hq}, \texttt{stag\_\allowbreak{}conj\_\allowbreak{}hq\_\allowbreak{}eq\_\allowbreak{}smul}, \texttt{corrQ\_\allowbreak{}hq\_\allowbreak{}neg}, \texttt{hq\_\allowbreak{}neg\_\allowbreak{}h}, \texttt{avgQ\_\allowbreak{}law\_\allowbreak{}neg\_\allowbreak{}h} \,(\texttt{Normalization.\allowbreak{}lean})}
\renewcommand{\stmtname}{\Cref{hyp:EC}}
\begin{stmtP}[EC]
There is a non-increasing $F:\N_0\to[0,1]$ with
\[
  M_p(F):=\sum_{d\ge0}(d+1)^pF(d)<\infty\quad\text{and}\quad M'_p(F):=\sum_{d\ge0}(d+1)^pF(d)^{1/2}<\infty
  \qquad\text{for }0\le p\le2,
\]
such that for all $N\ge1$ and all $x,y\in\{1,\dots,N\}$:
\begin{equation}\tag{\ref{eq:hypEC}}
  \E_q\,Q^{[1,N]}_q(x,y)\le F(|x-y|),
\end{equation}
where $Q^{[1,N]}_q$ is the eigenfunction correlator of $\hq$ on $\{1,\dots,N\}$ and $q$ is i.i.d.\
Bernoulli$(1/2)$.  (Theorem~\ref{thm:EC} proves this with $F(d)=\min\{1,Ce^{-\eta d}\}$.)
\end{stmtP}
\leanline{\texttt{corrQ}, \texttt{norm\_\allowbreak{}evo\_\allowbreak{}le\_\allowbreak{}corrQ}, \texttt{transProb\_\allowbreak{}le\_\allowbreak{}corrQ}, \texttt{EC} \,(\texttt{OneBody.\allowbreak{}lean})}
\renewcommand{\stmtname}{\Cref{lem:consequences}}
\begin{stmtP}[Consequences of correlator bounds]
Let $\Lambda$ be a finite interval and $k$ a real symmetric tridiagonal matrix on $\ell^2(\Lambda)$ whose
off-diagonal entries $k_{x,x+1}$ are all non-zero (e.g.\ $k=\hq$ or $k=A_\Lambda$).  Write $Q=Q_k$.
\begin{enumerate}[label=(\alph*)]
\item \emph{(Simple spectrum.)} $k$ has simple spectrum.  If $(E_\alpha,\varphi_\alpha)_\alpha$ is an eigenbasis with
$\varphi_\alpha$ real and orthonormal, then
\begin{equation}\tag{\ref{eq:Qsum}}
  Q(x,y)=\sum_\alpha|\varphi_\alpha(x)|\,|\varphi_\alpha(y)| ,
\end{equation}
and $Q(x,y)=Q(y,x)$, $Q(x,x)=1$, $0\le Q(x,y)\le1$.
\item \emph{(Functional calculus.)} For every bounded Borel $f$, $|f(k)_{xy}|\le\|f\|_\infty\,Q(x,y)$.  In
particular $\sup_{t\in\R}|(e^{-itk})_{xy}|\le Q(x,y)$, and $|P_{xy}|\le Q(x,y)$ for every spectral
projection $P=\one_S(k)$, e.g.\ the Fermi projections $\one_{(-\infty,\mu]}(k)$.
\item \emph{(Time-averaged transition probability.)} The limit
$R(x,y):=\lim_{T\to\infty}T^{-1}\int_0^T|(e^{-itk})_{xy}|^2\,dt$ exists and
\begin{equation}\tag{\ref{eq:Rdef}}
  R(x,y)=\sum_\alpha\varphi_\alpha(x)^2\varphi_\alpha(y)^2 ;
\end{equation}
moreover $R(x,y)=R(y,x)$, $\sum_{y\in\Lambda}R(x,y)=1$ and $R(x,y)\le Q(x,y)$.
\item \emph{(Annealed consequences.)} Assume Hypothesis~\ref{hyp:EC}, let $U_q(t):=e^{-it\hq}$ and let
$R_q$ be the quantity \eqref{eq:Rdef} for $k=\hq$.  Then for all $N\ge1$ and $x,y,z,x',y'\in[1,N]$:
(writing $Q=Q^{[1,N]}_q$)
\begin{align}
  &\E\sup_{t\in\R}|U_q(t)_{xy}|\le F(|x-y|),\qquad \E R_q(x,y)\le F(|x-y|),\qquad
   \E\,Q(x,y)^{1/2}\le F(|x-y|)^{1/2},\tag{\ref{eq:cons1}}\\
  &\E\sup_{t\in\R}|U_q(t)_{xz}|\,|U_q(t)_{yz}|\le\E\bigl[Q(x,z)\,Q(y,z)\bigr]\le F\bigl(\max\{|x-z|,|y-z|\}\bigr),\tag{\ref{eq:cons2}}\\
  &\E\bigl[Q(x,y)\,Q(x',y')\bigr]\le\bigl(F(|x-y|)\,F(|x'-y'|)\bigr)^{1/2},\tag{\ref{eq:cons3}}\\
  &\E\sup_{t\in\R}\sum_{x=1}^N(|x-y|+1)^p\,|U_q(t)_{xy}|^2\le\E\sum_{x=1}^N(|x-y|+1)^p\,Q(x,y)\le2M_p(F)\quad(0\le p\le2).\tag{\ref{eq:cons4}}
\end{align}
\end{enumerate}
\end{stmtP}
\leanline{\texttt{corrQ}, \texttt{norm\_\allowbreak{}evo\_\allowbreak{}le\_\allowbreak{}corrQ}, \texttt{transProb\_\allowbreak{}le\_\allowbreak{}corrQ}, \texttt{EC} \,(\texttt{OneBody.\allowbreak{}lean})}
\renewcommand{\stmtname}{\Cref{prop:nocritical}}
\begin{stmtP}[No critical energies]
Let $\lambda\neq0$ and $E\in\R$.  Let $G_E\subset SL(2,\R)$ be the closed subgroup generated by
$\{T_E(\lambda),T_E(-\lambda)\}$, the support of the law of $T_E(\lambda q_0)$.
\begin{enumerate}[label=(\roman*)]
\item $G_E$ is not compact.
\item $G_E$ is strongly irreducible: no finite non-empty set of lines in $\R^2$ is invariant under all of $G_E$.
\item The group generated by $T_E(\pm\lambda)$ is Zariski dense in $SL(2,\R)$.
\item $T_E(\lambda)T_E(-\lambda)\ne T_E(-\lambda)T_E(\lambda)$.  Regard the i.i.d.\ Bernoulli potential as
a random polymer model with polymers of length one.  Then no energy is critical in the sense of
[29, Def.~1]: there is no $E$ at which the polymer transfer matrices commute.
\end{enumerate}
Consequently $\gamma_\lambda(E)>0$ for every $E\in\R$.  This is Furstenberg's theorem
[36, Thm.~8.6], in the form [24, Thm.~2.1]: a probability measure on
$SL(2,\R)$ with $\int\log\|M\|\,d\nu<\infty$ whose closed generated group is non-compact and strongly
irreducible has positive Lyapunov exponent.  Moreover $E\mapsto\gamma_\lambda(E)$ is continuous
[24, Thm.~2.6], so
\begin{equation}\tag{\ref{eq:gstar}}
  \gstar(\lambda):=\min_{E\in\Ilam}\gamma_\lambda(E)>0 .
\end{equation}
\end{stmtP}
\leanline{\texttt{transferT\_\allowbreak{}mul\_\allowbreak{}inv}, \texttt{transferT\_\allowbreak{}commutator}, \texttt{transferT\_\allowbreak{}not\_\allowbreak{}commute}, \texttt{unipotent\_\allowbreak{}pow\_\allowbreak{}unbounded}, \texttt{no\_\allowbreak{}finite\_\allowbreak{}invariant\_\allowbreak{}lines}, \texttt{transferSL}, \texttt{transfer\_\allowbreak{}zariski\_\allowbreak{}dense} \,(\texttt{Furstenberg.\allowbreak{}lean}, \texttt{Zariski.\allowbreak{}lean})}
\renewcommand{\stmtname}{\Cref{lem:hypcheck}}
\begin{stmtP}[Hypotheses of the cited theorems]
For every $\lambda\ne0$, the operator $A_\Z=\Delta+\lambda\diag(q)$ on $\ell^2(\Z)$ satisfies:
\begin{enumerate}[label=(\roman*)]
\item the standing assumptions of [24, Def.~2.2]: an i.i.d.\ potential whose single-site law
has compact support $\mathcal A$ with $\#\mathcal A\ge2$.  Here $\mathcal A=\{-\lambda,\lambda\}$, and the
interval $\hat\Sigma=[-2-\max_{\alpha\in\mathcal A}|\alpha|,\,2+\max_{\alpha\in\mathcal A}|\alpha|]$ of
[24, Sect.~3] equals $\Ilam$;
\item assumptions (A)--(C) of [28] with $W=1$, $L_x\equiv1$, $V_x=\lambda q_x$;
\item the hypotheses of [27, Thm.~1.1]: $V_0$ is bounded and not concentrated at one point.
\end{enumerate}
\end{stmtP}
\leanline{verification of the hypotheses of the cited theorems in their own formalism, not formalized.}
\renewcommand{\stmtname}{\Cref{prop:LDT}}
\begin{stmtP}[Uniform large deviations; literature]
Let $\lambda\ne0$.  For every $\epsilon>0$ there are $C_\epsilon,c_\epsilon>0$, depending only on $\epsilon$
and $|\lambda|$, such that for all $m\ge1$, all $E\in\Ilam$ and all unit vectors $u\in\R^2$:
\begin{align}
  &\Prob\Bigl\{\Bigl|\tfrac1m\log\|\Phi_m(E)u\|-\gamma_\lambda(E)\Bigr|\ge\epsilon\Bigr\}\le C_\epsilon e^{-c_\epsilon m},\tag{\ref{eq:LDTvec}}\\
  &\Prob\Bigl\{\Bigl|\tfrac1m\log\|\Phi_m(E)\|-\gamma_\lambda(E)\Bigr|\ge\epsilon\Bigr\}\le C_\epsilon e^{-c_\epsilon m}.\tag{\ref{eq:LDTnorm}}
\end{align}
\end{stmtP}
\leanline{literature result (Furstenberg; uniform large deviations), not formalized: it enters the formalization as the hypothesis \texttt{LDT}.}
\renewcommand{\stmtname}{\Cref{prop:Zline}}
\begin{stmtP}[Ge--Zhao: annealed localization on the line]
Let $A_\Z=\Delta+V$ on $\ell^2(\Z)$ with $(V_x)$ i.i.d., $V_0$ bounded and not concentrated at one point,
and let $Q^\Z(x,y):=\sup\{|f(A_\Z)_{xy}|:f\text{ Borel},|f|\le1\}$.  Then for every $x\in\Z$
\begin{equation}\tag{\ref{eq:GeZhao}}
  \gamma^{\E}:=-\limsup_{y\to\pm\infty}\frac1{|y-x|}\log\E\,Q^\Z(x,y)>0
\end{equation}
[26], as stated in [27, Thm.~1.1].  In particular, for $V=\lambda q$ with
$\lambda\ne0$ and every $\eta'\in(0,\gamma^{\E})$ there is $C<\infty$ with
$\E\,Q^\Z(x,y)\le Ce^{-\eta'|x-y|}$ for all $x,y\in\Z$.
\end{stmtP}
\leanline{literature result (Ge--Zhao), not formalized; for Bernoulli potentials it also follows, given \texttt{LDT}, from the formalized Corollary~\ref{cor:ECinfvol}.}
\renewcommand{\stmtname}{\Cref{thm:EC}}
\begin{stmtP}[Hypothesis (EC) holds]
For every $\lambda\ne0$ there are constants $C<\infty$ and $\eta>0$, depending only on $|\lambda|$, such that:
\begin{enumerate}[label=(\roman*)]
\item for every finite interval $\Lambda\subset\Z$ and all $x,y\in\Lambda$,
\begin{equation}\tag{\ref{eq:ECmain}}
  \E\,Q^\Lambda(x,y)\le C\,e^{-\eta|x-y|};
\end{equation}
\item for all $x,y\in\Z$, with the supremum over all finite intervals $\Lambda\subset\Z$ containing $x$ and $y$,
\begin{equation}\tag{\ref{eq:ECsup}}
  \E\sup_{\Lambda\ni x,y}Q^\Lambda(x,y)\le C\,e^{-\eta|x-y|}.
\end{equation}
\end{enumerate}
In particular, by \cref{lem:normalization}, Hypothesis~\ref{hyp:EC} holds with $F(d)=\min\{1,Ce^{-\eta d}\}$.  For
this $F$, $M_p(F)\le C\sum_{d\ge0}(d+1)^pe^{-\eta d}$ and $M'_p(F)\le C^{1/2}\sum_{d\ge0}(d+1)^pe^{-\eta d/2}$.
\end{stmtP}
\leanline{\texttt{ec\_\allowbreak{}of\_\allowbreak{}ldt}, \texttt{corrQ\_\allowbreak{}decay\_\allowbreak{}of\_\allowbreak{}ldt}, \texttt{poisson}, \texttt{box\_\allowbreak{}right}, \texttt{box\_\allowbreak{}left}, \texttt{dirL\_\allowbreak{}eq\_\allowbreak{}dirR}, \texttt{sl2\_\allowbreak{}entry\_\allowbreak{}ineq}, \texttt{eigvec\_\allowbreak{}recursion}, \texttt{abs\_\allowbreak{}eigenvalue\_\allowbreak{}hq\_\allowbreak{}le}, \texttt{decay\_\allowbreak{}of\_\allowbreak{}not\_\allowbreak{}res}, \texttt{exists\_\allowbreak{}resEnd}, \texttt{card\_\allowbreak{}resEnds}, \texttt{LDT}, \texttt{entry\_\allowbreak{}ldt}, \texttt{single\_\allowbreak{}box}, \texttt{two\_\allowbreak{}box} \,(\texttt{Transfer.\allowbreak{}lean}, \texttt{Resonance.\allowbreak{}lean}, \texttt{LDTEntry.\allowbreak{}lean}, \texttt{ECProof.\allowbreak{}lean});\\ \texttt{ec\_\allowbreak{}sup\_\allowbreak{}of\_\allowbreak{}ldt}, \texttt{corrQwin}, \texttt{corrQsup} \,(\texttt{ECSup.\allowbreak{}lean})}
\renewcommand{\stmtname}{\Cref{cor:ECinfvol}}
\begin{stmtP}[Infinite volume: half-line and line]
Let $\lambda\neq0$, $X\in\{\N,\Z\}$ with $\N=\{1,2,\dots\}$, and $A_X:=\Delta_X+\lambda\diag(q)$ on
$\ell^2(X)$, with the Dirichlet condition $\psi(0)=0$ if $X=\N$.  Let
$Q^X(x,y):=\sup\{|f(A_X)_{xy}|:f\text{ Borel},|f|\le1\}$, and let $\Lambda_L:=[1,L]$ if $X=\N$ and
$\Lambda_L:=[-L,L]$ if $X=\Z$.  Then for all $x,y\in X$
\begin{equation}\tag{\ref{eq:ECinf}}
  Q^X(x,y)\le\liminf_{L\to\infty}Q^{\Lambda_L}(x,y)\ \ \text{pointwise, hence}\ \
  \E\,Q^X(x,y)\le\liminf_{L\to\infty}\E\,Q^{\Lambda_L}(x,y)\le Ce^{-\eta|x-y|},
\end{equation}
with $C,\eta$ from Theorem~\ref{thm:EC}.  In particular
$\E\sup_t|\langle\delta_x,e^{-itA_X}\delta_y\rangle|\le Ce^{-\eta|x-y|}$ on the half-line with Dirichlet
boundary condition.  For $X=\Z$ this gives an alternative proof of \cref{prop:Zline} for Bernoulli
potentials, with $\gamma^\E\ge\eta$.
\end{stmtP}
\leanline{\texttt{tendsto\_\allowbreak{}evoZ}, \texttt{lintegral\_\allowbreak{}iSup\_\allowbreak{}UZ\_\allowbreak{}le}, \texttt{infvol\_\allowbreak{}propagator\_\allowbreak{}of\_\allowbreak{}ldt}, \texttt{tendsto\_\allowbreak{}density\_\allowbreak{}one}, \texttt{moments\_\allowbreak{}UZ}, \texttt{abs\_\allowbreak{}transportM\_\allowbreak{}le}, \texttt{transport\_\allowbreak{}tendsto} \,(\texttt{InfVol.\allowbreak{}lean}, \texttt{InfVolTransport.\allowbreak{}lean});\\ \texttt{IsAZ}, \texttt{exists\_\allowbreak{}isAZ}, \texttt{corrQZ}, \texttt{lintegral\_\allowbreak{}corrQZ\_\allowbreak{}le}, \texttt{IsAN}, \texttt{exists\_\allowbreak{}isAN}, \texttt{corrQN}, \texttt{lintegral\_\allowbreak{}corrQN\_\allowbreak{}le} \,(\texttt{ECInfVolQ.\allowbreak{}lean}, \texttt{ECInfVolQN.\allowbreak{}lean})}
\renewcommand{\stmtname}{\Cref{lem:green}}
\begin{stmtP}[Green functions through one transfer-matrix product]
Define $u_\pm:[a-1,b+1]\to\R$ by $u_-(a-1)=0$, $u_-(a)=1$, $u_+(b+1)=0$, $u_+(b)=1$, and
\begin{equation}\tag{\ref{eq:rec}}
  u(z+1)+u(z-1)+v_zu(z)=Eu(z),\qquad z\in[a,b]
\end{equation}
for both $u=u_\pm$.
\begin{enumerate}[label=(\alph*)]
\item $u_-(z)=\det(E-A_{[a,z-1]})$ and $u_+(z)=\det(E-A_{[z+1,b]})$ for $z\in[a,b+1]$, resp.\
$z\in[a-1,b]$.  The Wronskian $\Wr:=u_-(z+1)u_+(z)-u_-(z)u_+(z+1)$ does not depend on $z\in[a-1,b]$,
and $\Wr=\det(E-A_B)$.
\item If $E\notin\sigma(A_B)$, then $G_E[A_B](z,z')=G_E[A_B](z',z)=-u_-(z)u_+(z')/\Wr$ for $a\le z\le z'\le b$.
\item \emph{(Right end.)} Let $x\in[a,b-1]$, $M:=T_{x+1}T_{x+2}\cdots T_b$ (a product of $b-x$ matrices),
$Y:=(u_-(x+1),u_-(x))^\top$ and $w:=DY/\|Y\|$.  Then $Y\ne0$, the unit vector $w$ is a function of
$(v_z)_{a\le z\le x}$ only, and
\begin{equation}\tag{\ref{eq:greenR}}
  \Wr=\|Y\|\,\langle w,M\eone\rangle,\qquad
  \bigl|\det(E-A_{[a,x-1]})\bigr|\,\bigl|\langle w,M\eone\rangle\bigr|\le\bigl|\det(E-A_B)\bigr| .
\end{equation}
In particular $E\in\sigma(A_B)$ if and only if $\langle w,M\eone\rangle=0$, and otherwise
$|G_E[A_B](x,b)|\le|\langle w,M\eone\rangle|^{-1}$.
\item \emph{(Left end.)} Let $x\in[a+1,b]$, $M':=T_{x-1}T_{x-2}\cdots T_a$ (a product of $x-a$ matrices),
$Y':=(u_+(x-1),u_+(x))^\top$ and $w':=DY'/\|Y'\|$.  Then $Y'\neq0$, $w'$ is a function of $(v_z)_{x\le z\le b}$
only, and
\begin{equation}\tag{\ref{eq:greenL}}
  \Wr=\|Y'\|\,\langle w',M'\eone\rangle,\qquad
  \bigl|\det(E-A_{[x+1,b]})\bigr|\,\bigl|\langle w',M'\eone\rangle\bigr|\le\bigl|\det(E-A_B)\bigr| .
\end{equation}
\end{enumerate}
\end{stmtP}
\leanline{\texttt{ec\_\allowbreak{}of\_\allowbreak{}ldt}, \texttt{corrQ\_\allowbreak{}decay\_\allowbreak{}of\_\allowbreak{}ldt}, \texttt{poisson}, \texttt{box\_\allowbreak{}right}, \texttt{box\_\allowbreak{}left}, \texttt{dirL\_\allowbreak{}eq\_\allowbreak{}dirR}, \texttt{sl2\_\allowbreak{}entry\_\allowbreak{}ineq}, \texttt{eigvec\_\allowbreak{}recursion}, \texttt{abs\_\allowbreak{}eigenvalue\_\allowbreak{}hq\_\allowbreak{}le}, \texttt{decay\_\allowbreak{}of\_\allowbreak{}not\_\allowbreak{}res}, \texttt{exists\_\allowbreak{}resEnd}, \texttt{card\_\allowbreak{}resEnds}, \texttt{LDT}, \texttt{entry\_\allowbreak{}ldt}, \texttt{single\_\allowbreak{}box}, \texttt{two\_\allowbreak{}box} \,(\texttt{Transfer.\allowbreak{}lean}, \texttt{Resonance.\allowbreak{}lean}, \texttt{LDTEntry.\allowbreak{}lean}, \texttt{ECProof.\allowbreak{}lean})}
\renewcommand{\stmtname}{\Cref{lem:entry}}
\begin{stmtP}[Large deviations for entries in arbitrary directions]
Let $\lambda\ne0$ and $\gstar=\gstar(\lambda)$ as in \eqref{eq:gstar}.  There are $C_1,c_1>0$ and
$m_0\in\N$, depending only on $|\lambda|$, such that the following holds for every $m\ge m_0$, every
$E\in\Ilam$, all unit vectors $e,w\in\R^2$, and every product $M$ of $m$ matrices $T_E(\lambda q_z)$
over $m$ distinct sites $z$ (in any order):
\begin{equation}\tag{\ref{eq:entryLDT}}
  \Prob\bigl\{|\langle w,Me\rangle|\le e^{\gstar m/2}\bigr\}\le C_1e^{-c_1m}.
\end{equation}
\end{stmtP}
\leanline{\texttt{ec\_\allowbreak{}of\_\allowbreak{}ldt}, \texttt{corrQ\_\allowbreak{}decay\_\allowbreak{}of\_\allowbreak{}ldt}, \texttt{poisson}, \texttt{box\_\allowbreak{}right}, \texttt{box\_\allowbreak{}left}, \texttt{dirL\_\allowbreak{}eq\_\allowbreak{}dirR}, \texttt{sl2\_\allowbreak{}entry\_\allowbreak{}ineq}, \texttt{eigvec\_\allowbreak{}recursion}, \texttt{abs\_\allowbreak{}eigenvalue\_\allowbreak{}hq\_\allowbreak{}le}, \texttt{decay\_\allowbreak{}of\_\allowbreak{}not\_\allowbreak{}res}, \texttt{exists\_\allowbreak{}resEnd}, \texttt{card\_\allowbreak{}resEnds}, \texttt{LDT}, \texttt{entry\_\allowbreak{}ldt}, \texttt{single\_\allowbreak{}box}, \texttt{two\_\allowbreak{}box} \,(\texttt{Transfer.\allowbreak{}lean}, \texttt{Resonance.\allowbreak{}lean}, \texttt{LDTEntry.\allowbreak{}lean}, \texttt{ECProof.\allowbreak{}lean})}
\renewcommand{\stmtname}{\Cref{lem:badset}}
\begin{stmtP}[Structure of the resonant set]
\leavevmode
\begin{enumerate}[label=(\alph*)]
\item $\Bad$ is closed and contains $\sigma(A_B)\cap\Ilam$.  For $E\in\Ilam\setminus\Bad$: $E\notin\sigma(A_B)$ and
$|G_E[A_B](x,e)|<e^{-\kappa n}$ for every $e\in\Ends$.
\item There is a set $\mathcal Z\subset\Ilam$ with $\#\mathcal Z\le8n+6$, determined by $(v_z)_{z\in B}$, such
that every connected component of $\Bad$ is a closed interval (possibly a point) with both endpoints in
$\mathcal Z$.
\end{enumerate}
\end{stmtP}
\leanline{\texttt{ec\_\allowbreak{}of\_\allowbreak{}ldt}, \texttt{corrQ\_\allowbreak{}decay\_\allowbreak{}of\_\allowbreak{}ldt}, \texttt{poisson}, \texttt{box\_\allowbreak{}right}, \texttt{box\_\allowbreak{}left}, \texttt{dirL\_\allowbreak{}eq\_\allowbreak{}dirR}, \texttt{sl2\_\allowbreak{}entry\_\allowbreak{}ineq}, \texttt{eigvec\_\allowbreak{}recursion}, \texttt{abs\_\allowbreak{}eigenvalue\_\allowbreak{}hq\_\allowbreak{}le}, \texttt{decay\_\allowbreak{}of\_\allowbreak{}not\_\allowbreak{}res}, \texttt{exists\_\allowbreak{}resEnd}, \texttt{card\_\allowbreak{}resEnds}, \texttt{LDT}, \texttt{entry\_\allowbreak{}ldt}, \texttt{single\_\allowbreak{}box}, \texttt{two\_\allowbreak{}box} \,(\texttt{Transfer.\allowbreak{}lean}, \texttt{Resonance.\allowbreak{}lean}, \texttt{LDTEntry.\allowbreak{}lean}, \texttt{ECProof.\allowbreak{}lean})}
\renewcommand{\stmtname}{\Cref{lem:double}}
\begin{stmtP}[Two-point resonance elimination]
Let $\lambda\ne0$, and let $C_1,c_1,m_0$ be as in \cref{lem:entry}.  Let $\Lambda$ be a finite interval,
$x<y$ in $\Lambda$, $d:=y-x$, and $n:=\lfloor(d-1)/2\rfloor$; assume $n\ge m_0$.  Let
\[
  B_x:=[x-n,x+n]\cap\Lambda,\qquad B_y:=[y-n,y+n]\cap\Lambda,
\]
\[
  \Ends_x:=\{x+n\}\cup\bigl(\{x-n\}\text{ if }x-n-1\in\Lambda\bigr),\qquad
  \Ends_y:=\{y-n\}\cup\bigl(\{y+n\}\text{ if }y+n+1\in\Lambda\bigr),
\]
and let $\Bad_x,\Bad_y$ and $\mathcal Z_x,\mathcal Z_y$ be the corresponding resonant sets and endpoint sets
(\cref{lem:badset}).  Then $B_x\cap B_y=\emptyset$; the elements of $\Ends_x$ and $\Ends_y$ are ends of
$B_x$, resp.\ $B_y$, at distance $n$; and
\begin{gather}
  \sup_{E\in\Ilam}\Prob\{E\in\Bad_x\}\le2C_1e^{-c_1n},\qquad
  \sup_{E\in\Ilam}\Prob\{E\in\Bad_y\}\le2C_1e^{-c_1n},\tag{\ref{eq:single}}\\
  \Prob\{\Bad_x\cap\Bad_y\neq\emptyset\}\le4(8n+6)\,C_1e^{-c_1n}.\tag{\ref{eq:double}}
\end{gather}
\end{stmtP}
\leanline{\texttt{ec\_\allowbreak{}of\_\allowbreak{}ldt}, \texttt{corrQ\_\allowbreak{}decay\_\allowbreak{}of\_\allowbreak{}ldt}, \texttt{poisson}, \texttt{box\_\allowbreak{}right}, \texttt{box\_\allowbreak{}left}, \texttt{dirL\_\allowbreak{}eq\_\allowbreak{}dirR}, \texttt{sl2\_\allowbreak{}entry\_\allowbreak{}ineq}, \texttt{eigvec\_\allowbreak{}recursion}, \texttt{abs\_\allowbreak{}eigenvalue\_\allowbreak{}hq\_\allowbreak{}le}, \texttt{decay\_\allowbreak{}of\_\allowbreak{}not\_\allowbreak{}res}, \texttt{exists\_\allowbreak{}resEnd}, \texttt{card\_\allowbreak{}resEnds}, \texttt{LDT}, \texttt{entry\_\allowbreak{}ldt}, \texttt{single\_\allowbreak{}box}, \texttt{two\_\allowbreak{}box} \,(\texttt{Transfer.\allowbreak{}lean}, \texttt{Resonance.\allowbreak{}lean}, \texttt{LDTEntry.\allowbreak{}lean}, \texttt{ECProof.\allowbreak{}lean})}
\subsection*{\S\,\ref{sec:main}\quad Main results: disorder-free localization}
\addcontentsline{toc}{subsection}{\S\,\ref{sec:main}\quad Main results: disorder-free localization}
\renewcommand{\stmtname}{\Cref{lem:sectornorm}}
\begin{stmtP}[Sector-diagonal operators]
Let $X$ be a gauge-invariant operator on $\cH_N$.
\begin{enumerate}[label=(\alph*)]
\item $X=\sum_qP_qXP_q$ and $\norm{X}=\max_q\norm{P_qXP_q}$.
\item Let $\rho$ be a density matrix on $\cH_N$ and $w_q:=\Tr(\rho P_q)$ its sector weights.  Then
$|\Tr(\rho X)|\le\sum_qw_q\norm{P_qXP_q}$.  One has $w_q=2^{-N}$, and hence
$|\Tr(\rho X)|\le\E_q\norm{P_qXP_q}$, for $\rho=4^{-N}\one$ and for $\rho=\ket{\Psi_0}\bra{\Psi_0}$ with
$\Psi_0=\ket{\uparrow\cdots\uparrow}_z\otimes\psi$, where $\psi\in\cF_N$ is any unit vector.
\item If $X=\sum_qx_qP_q$ with $x_q\in\C$, then
$\langle\Psi_0,|X|\Psi_0\rangle=4^{-N}\Tr|X|=\E_q|x_q|$ for every $\Psi_0$ as in (b).
\end{enumerate}
\end{stmtP}
\leanline{\texttt{gaugeInv\_\allowbreak{}eq\_\allowbreak{}sum}, \texttt{norm\_\allowbreak{}gaugeInv}, \texttt{norm\_\allowbreak{}trace\_\allowbreak{}le}, \texttt{trace\_\allowbreak{}mixed\_\allowbreak{}gaussProj}, \texttt{trace\_\allowbreak{}initZ\_\allowbreak{}gaussProj}, \texttt{norm\_\allowbreak{}trace\_\allowbreak{}mixed\_\allowbreak{}le}, \texttt{norm\_\allowbreak{}expect\_\allowbreak{}initZ\_\allowbreak{}le}, \texttt{sectorDiag\_\allowbreak{}abs}, \texttt{expect\_\allowbreak{}sectorDiag\_\allowbreak{}abs}, \texttt{trace\_\allowbreak{}sectorDiag\_\allowbreak{}abs} \,(\texttt{SectorNorm.\allowbreak{}lean})}
\renewcommand{\stmtname}{\Cref{lem:quasifree}}
\begin{stmtP}[Quasi-free density matrices]
Let $0\le C\le\one$ on $\ell^2(A)$.
\begin{enumerate}[label=(\alph*)]
\item With $N_A:=\sum_{x\in A}c_x^\dagger c_x$ and the convention that the $j=0$ term is $\Gamma(C)$,
\begin{equation}\tag{\ref{eq:main-rhoformula}}
  \rhoqf{C}=\sum_{j=0}^{|A|}\frac{(-1)^j}{j!}\sum_{x_1,\dots,x_j\in A}c_{x_1}\cdots c_{x_j}\,\Gamma(C)\,
  c^\dagger_{x_j}\cdots c^\dagger_{x_1}.
\end{equation}
In particular $\rhoqf{C}$ does not depend on the choice of eigenbasis, its matrix elements in the
occupation basis are polynomials in the entries of $C$, and $C\mapsto\rhoqf{C}$ is uniformly continuous
on $\{0\le C\le\one\}$.
\item $\rhoqf{C}$ is a density matrix with $S(\rhoqf{C})=\sum_i\bent(\nu_i)$.  Under $\rhoqf{C}$, $N_A$ has
the law of $\sum_i\xi_i$ with independent $\xi_i\sim\mathrm{Bernoulli}(\nu_i)$; its generating function is
$\Tr(\rhoqf{C}z^{N_A})=\det(\one+(z-1)C)$, and its variance is $\Tr(C-C^2)$.
\end{enumerate}
\end{stmtP}
\leanline{\texttt{corrA\_\allowbreak{}isHermitian}, \texttt{eig\_\allowbreak{}corrA\_\allowbreak{}mem}, \texttt{vonNeumann\_\allowbreak{}rdm\_\allowbreak{}slater}, \texttt{numDist\_\allowbreak{}slater}, \texttt{genfun\_\allowbreak{}slater}, \texttt{sum\_\allowbreak{}sqrt\_\allowbreak{}eig\_\allowbreak{}le} \,(\texttt{SlaterExact.\allowbreak{}lean});\\ \texttt{corrR\_\allowbreak{}isHermitian}, \texttt{eig\_\allowbreak{}corrR\_\allowbreak{}mem}, \texttt{genfun\_\allowbreak{}slater\_\allowbreak{}gen}, \texttt{numDist\_\allowbreak{}slater\_\allowbreak{}gen} \,(\texttt{SlaterGeneral.\allowbreak{}lean});\\ \texttt{secondQ}, \texttt{qfState}, \texttt{qfState\_\allowbreak{}eq\_\allowbreak{}prod}, \texttt{continuous\_\allowbreak{}qfState}, \texttt{uniformContinuousOn\_\allowbreak{}qfState}, \texttt{qfState\_\allowbreak{}posSemidef}, \texttt{trace\_\allowbreak{}qfState}, \texttt{vonNeumann\_\allowbreak{}qfState}, \texttt{trace\_\allowbreak{}qfState\_\allowbreak{}zN}, \texttt{rdmF\_\allowbreak{}eq\_\allowbreak{}qfState} \,(\texttt{QuasiFree.\allowbreak{}lean})}
\renewcommand{\stmtname}{\Cref{lem:slater}}
\begin{stmtP}[Reduced states of Slater determinants]
Let $\Lambda$ be a finite set of modes, $W\subset\ell^2(\Lambda)$ a subspace of dimension $M$ with orthogonal
projection $K$, and $\psi_W:=c^\dagger(w_1)\cdots c^\dagger(w_M)\ket0$ for an orthonormal basis $(w_j)$ of $W$
(unique up to a phase).  Let $A\subset\Lambda$, $B:=\Lambda\setminus A$, and identify
$\cF_\Lambda\cong\cF_A\otimes\cF_B$ by $c(a)\mapsto c(a)\otimes\one$, $c(b)\mapsto(-1)^{N_A}\otimes c(b)$ for
$a\in\ell^2(A)$, $b\in\ell^2(B)$, $\ket0\mapsto\ket0\otimes\ket0$.  (For $A=\{1,\dots,\ell\}$ this is the
identification of \S\ref{sec:model:cut}.)  Under it, the CAR algebra of the modes in $A$ is
$\mathcal B(\cF_A)\otimes\one$, and we let $\rho_A:=\Tr_{\cF_B}\ket{\psi_W}\bra{\psi_W}$.  Let
$C_A:=P_AKP_A$, regarded as an operator on $\ell^2(A)$, with eigenvalues $\nu_1,\dots,\nu_{|A|}$, and
$T:=P_BKP_A$.  Then:
\begin{enumerate}[label=(\alph*)]
\item $\rho_A=\rhoqf{C_A}$.  Hence $S(\rho_A)=\sum_i\bent(\nu_i)$, and the law of $N_A$ under $\psi_W$ is the
Poisson-binomial law of \cref{lem:quasifree}(b).
\item $C_A-C_A^2=T^*T$.  Hence
$\Var_{\psi_W}(N_A)=\norm{T}_2^2=\sum_{x\in A,\,y\in B}|K_{yx}|^2$ and
$\sum_i\sqrt{\nu_i(1-\nu_i)}=\norm{T}_1\le\sum_{x\in A,\,y\in B}|K_{yx}|$.
\end{enumerate}
\end{stmtP}
\leanline{\texttt{corrA\_\allowbreak{}isHermitian}, \texttt{eig\_\allowbreak{}corrA\_\allowbreak{}mem}, \texttt{vonNeumann\_\allowbreak{}rdm\_\allowbreak{}slater}, \texttt{numDist\_\allowbreak{}slater}, \texttt{genfun\_\allowbreak{}slater}, \texttt{sum\_\allowbreak{}sqrt\_\allowbreak{}eig\_\allowbreak{}le} \,(\texttt{SlaterExact.\allowbreak{}lean});\\ \texttt{corrR\_\allowbreak{}isHermitian}, \texttt{eig\_\allowbreak{}corrR\_\allowbreak{}mem}, \texttt{genfun\_\allowbreak{}slater\_\allowbreak{}gen}, \texttt{numDist\_\allowbreak{}slater\_\allowbreak{}gen} \,(\texttt{SlaterGeneral.\allowbreak{}lean});\\ \texttt{rdmF}, \texttt{reorderSign}, \texttt{vonNeumann\_\allowbreak{}rdmF\_\allowbreak{}slater}, \texttt{vonNeumann\_\allowbreak{}rdmF\_\allowbreak{}le}, \texttt{vonNeumann\_\allowbreak{}rdmF\_\allowbreak{}segment}, \texttt{lgt\_\allowbreak{}areaLaw\_\allowbreak{}bulk} \,(\texttt{SlaterBulk.\allowbreak{}lean});\\ \texttt{secondQ}, \texttt{qfState}, \texttt{qfState\_\allowbreak{}eq\_\allowbreak{}prod}, \texttt{continuous\_\allowbreak{}qfState}, \texttt{uniformContinuousOn\_\allowbreak{}qfState}, \texttt{qfState\_\allowbreak{}posSemidef}, \texttt{trace\_\allowbreak{}qfState}, \texttt{vonNeumann\_\allowbreak{}qfState}, \texttt{trace\_\allowbreak{}qfState\_\allowbreak{}zN}, \texttt{rdmF\_\allowbreak{}eq\_\allowbreak{}qfState} \,(\texttt{QuasiFree.\allowbreak{}lean})}
\renewcommand{\stmtname}{\Cref{lem:binent}}
\begin{stmtP}[Binary entropy]
For all $\nu\in[0,1]$: $\bent(\nu)\le2\ln2\,\sqrt{\nu(1-\nu)}$, with equality iff $\nu\in\{0,\tfrac12,1\}$.
\end{stmtP}
\leanline{\texttt{binEntropy\_\allowbreak{}le\_\allowbreak{}two\_\allowbreak{}log\_\allowbreak{}two\_\allowbreak{}sqrt}, \texttt{binEntropy\_\allowbreak{}eq\_\allowbreak{}two\_\allowbreak{}log\_\allowbreak{}two\_\allowbreak{}sqrt\_\allowbreak{}iff} \,(\texttt{BinEntropy.\allowbreak{}lean})}
\renewcommand{\stmtname}{\Cref{lem:annealed}}
\begin{stmtP}[Annealed one-body bounds]
Assume Hypothesis~\ref{hyp:EC}.  Let $\psi$ be a Fock state with occupied set $S$ (arbitrary), and
$x\neq y$ in $[1,N]$, $d:=|x-y|$.
\begin{enumerate}[label=(\roman*)]
\item $\displaystyle\E\sup_{t\in\R}|K_q(t)_{xy}|\le\E\sum_{z=1}^NQ_q(x,z)Q_q(y,z)\le2\sum_{k\ge\lceil d/2\rceil}F(k)$.
\item For $1\le m<N$:
\begin{gather*}
  \E\sup_t\sum_{x\le m<y}|K_q(t)_{xy}|\le\ThF,\qquad
  \E\sum_{x\le m<z}Q_q(x,z)\le\SigF,\\
  \E\Bigl(\sum_{x\le m<z}Q_q(x,z)\Bigr)^2\le\SigpF^2 .
\end{gather*}
\item For an interval $B\subset[1,N]$ the same bounds hold for the sums over pairs $x\in B$,
$y\in[1,N]\setminus B$, with the constants $2\ThF$, $2\SigF$ and $4\SigpF^2$.
\end{enumerate}
\end{stmtP}
\leanline{\texttt{annealed\_\allowbreak{}K}, \texttt{annealed\_\allowbreak{}cut}, \texttt{vonNeumann\_\allowbreak{}slater\_\allowbreak{}sharp}, \texttt{lgt\_\allowbreak{}areaLaw\_\allowbreak{}sharp}, \texttt{thetaF\_\allowbreak{}le}, \texttt{lgt\_\allowbreak{}areaLaw\_\allowbreak{}sharp\_\allowbreak{}of\_\allowbreak{}ldt} \,(\texttt{AreaLawSharp.\allowbreak{}lean});\\ \texttt{annealed\_\allowbreak{}cut\_\allowbreak{}Q}, \texttt{annealed\_\allowbreak{}cut\_\allowbreak{}Q\_\allowbreak{}sq}, \texttt{annealed\_\allowbreak{}ivl\_\allowbreak{}K}, \texttt{annealed\_\allowbreak{}ivl\_\allowbreak{}Q}, \texttt{annealed\_\allowbreak{}ivl\_\allowbreak{}Q\_\allowbreak{}sq} \,(\texttt{AnnealedInterval.\allowbreak{}lean})}
\renewcommand{\stmtname}{\Cref{lem:propagation}}
\begin{stmtP}[Finite propagation speed]
Let $\Lambda\subset\Lambda'\subset\Z$ be intervals ($\Lambda'=\Z$ allowed), $v:\Lambda'\to\R$ with $|v|\le V_0$,
$k_{\Lambda}:=-J\Delta_\Lambda+\diag(v|_\Lambda)$ on $\ell^2(\Lambda)$, likewise $k_{\Lambda'}$, and $c:=|t|(2J+V_0)$.  For
$x,y\in\Lambda$:
\begin{enumerate}[label=(\alph*)]
\item $|(e^{-itk_\Lambda})_{xy}|\le e^{c}\,c^{|x-y|}/|x-y|!$;
\item $|(e^{-itk_\Lambda})_{xy}-(e^{-itk_{\Lambda'}})_{xy}|\le2e^{c}\,c^{n_0}/n_0!$, where
$n_0:=\dist(x,\Lambda'\setminus\Lambda)+\dist(y,\Lambda'\setminus\Lambda)$ (the left side is $0$ if $\Lambda=\Lambda'$).
\end{enumerate}
\end{stmtP}
\leanline{\texttt{tendsto\_\allowbreak{}evoZ}, \texttt{lintegral\_\allowbreak{}iSup\_\allowbreak{}UZ\_\allowbreak{}le}, \texttt{infvol\_\allowbreak{}propagator\_\allowbreak{}of\_\allowbreak{}ldt}, \texttt{tendsto\_\allowbreak{}density\_\allowbreak{}one}, \texttt{moments\_\allowbreak{}UZ}, \texttt{abs\_\allowbreak{}transportM\_\allowbreak{}le}, \texttt{transport\_\allowbreak{}tendsto} \,(\texttt{InfVol.\allowbreak{}lean}, \texttt{InfVolTransport.\allowbreak{}lean})}
\renewcommand{\stmtname}{\Cref{thm:transport}}
\begin{stmtP}[Dressed-fermion localization and bounded transport]
Assume Hypothesis~\ref{hyp:EC}, and let $\Psi_0=\ket{\uparrow\cdots\uparrow}_z\otimes\psi$,
$\Psi(t)=e^{-itH}\Psi_0$.  For all $N\ge2$:
\begin{enumerate}[label=(\alph*)]
\item \emph{(Propagator.)}  For all $x,y\in[1,N]$ and $t\in\R$,
\begin{equation}\tag{\ref{eq:main-anticomm}}
  \{\cdr_x(t),\cdr_y^\dagger\}=\sum_qU_q(t)_{xy}\,P_q .
\end{equation}
Thus the dressed-fermion propagator is the scalar $U_q(t)_{xy}$ in each sector, and its sector-averaged
modulus is localized uniformly in time:
\begin{equation}\tag{\ref{eq:main-prop}}
  \E_q\sup_{t\in\R}\norm{P_q\{\cdr_x(t),\cdr_y^\dagger\}P_q}=\E_q\sup_{t\in\R}|U_q(t)_{xy}|\le F(|x-y|).
\end{equation}
In particular, $\sup_t\langle\Psi_0,|\{\cdr_x(t),\cdr_y^\dagger\}|\,\Psi_0\rangle\le F(|x-y|)$ for every
matter state $\psi$.
\item \emph{(One particle: moments.)}  If $\psi=f_y^\dagger\ket0$, then for $0\le p\le2$,
\[
  \sup_{t\in\R}\sum_{x=1}^N(|x-y|+1)^p\langle\Psi(t),n_x\Psi(t)\rangle
  \le\E_q\sup_{t\in\R}\sum_{x=1}^N(|x-y|+1)^p|U_q(t)_{xy}|^2\le2M_p(F).
\]
\item \emph{(Transport across any cut.)}  Let $\psi=\ket n$ be a Fock state with occupied set $S$, let
$1\le m<N$ and $s_m:=|S\cap[1,m]|$.  Then for every $q$ and $t$,
$|\langle\psi_q(t),N_{\le m}\psi_q(t)\rangle-s_m|\le\sum_{x\le m<z}|U_q(t)_{xz}|^2$, and
\begin{gather}
  \E_q\sup_t\bigl|\langle\psi_q(t),N_{\le m}\psi_q(t)\rangle-s_m\bigr|\le\SigF\le M_1(F),\qquad
  \sup_t\bigl|\langle\Psi(t),N_{\le m}\Psi(t)\rangle-s_m\bigr|\le\SigF,\tag{\ref{eq:main-transport}}\\
  \sup_t\bigl\langle\Psi(t),(N_{\le m}-s_m)^2\Psi(t)\bigr\rangle\le\ThF+\SigpF^2\le4M_2(F)+M'_1(F)^2.
  \tag{\ref{eq:main-transport2}}
\end{gather}
In particular $\Prob_{\Psi(t)}(|N_{\le m}-s_m|\ge k)\le(4M_2(F)+M'_1(F)^2)/k^2$ for all $t$ and $k>0$.
\item \emph{(Dressed correlations.)}  For Fock $\psi$ and $x\neq y$,
$\sup_t|\langle\Psi(t),\cdr_x^\dagger\cdr_y\Psi(t)\rangle|\le2\sum_{k\ge\lceil|x-y|/2\rceil}F(k)$; here
$\cdr_x^\dagger\cdr_y=f_x^\dagger(\prod_{x\le k<y}\sigma^z_{b_k})f_y$ for $x<y$ (\cref{thm:duality}(d)).
\end{enumerate}
All bounds are uniform in $N$, in the position of the cut and in $t$.
\end{stmtP}
\leanline{\texttt{lgt\_\allowbreak{}density}, \texttt{lgt\_\allowbreak{}transport} \,(\texttt{Main.\allowbreak{}lean});\\ \texttt{dressed\_\allowbreak{}anticomm\_\allowbreak{}t}, \texttt{norm\_\allowbreak{}sector\_\allowbreak{}anticomm}, \texttt{avg\_\allowbreak{}sup\_\allowbreak{}anticomm\_\allowbreak{}le}, \texttt{expect\_\allowbreak{}abs\_\allowbreak{}anticomm\_\allowbreak{}le}, \texttt{one\_\allowbreak{}particle\_\allowbreak{}moments}, \texttt{sector\_\allowbreak{}transport\_\allowbreak{}le}, \texttt{avg\_\allowbreak{}sup\_\allowbreak{}transport\_\allowbreak{}le}, \texttt{lgt\_\allowbreak{}transport\_\allowbreak{}sharp}, \texttt{lgt\_\allowbreak{}transport\_\allowbreak{}sq}, \texttt{dressed\_\allowbreak{}corr\_\allowbreak{}le}, \texttt{lgt\_\allowbreak{}numberEntropy\_\allowbreak{}sharp}, \texttt{avg\_\allowbreak{}sup\_\allowbreak{}sector\_\allowbreak{}numberEntropy\_\allowbreak{}le} \,(\texttt{TransportSharp.\allowbreak{}lean})}
\renewcommand{\stmtname}{\Cref{prop:LR}}
\begin{stmtP}[Sector-averaged zero-velocity Lieb--Robinson bound]
Let $X,Y\subset[1,N]$ and $B\in\BAlg^+(Y)$.
\begin{enumerate}[label=(\alph*)]
\item \emph{(Sector-wise bound; no hypothesis.)}  For a monomial $A=\cdr^{\sharp}_{x_1}\cdots\cdr^{\sharp}_{x_k}$
($x_i\in X$, each $\cdr^\sharp$ equal to $\cdr$ or $\cdr^\dagger$), and for every $q$ and $t$,
\begin{equation}\tag{\ref{eq:main-LRsector}}
  \norm{P_q[A(t),B]}\le2\norm B\sum_{i=1}^k\sum_{y\in Y}|U_q(t)_{x_iy}| .
\end{equation}
For general $A\in\Alg(X)$, $\norm{P_q[A(t),B]}\le8|A|_1\norm B\sum_{x\in X,y\in Y}|U_q(t)_{xy}|$.
\item \emph{(Sector-averaged zero-velocity bound.)}  Under Hypothesis~\ref{hyp:EC},
\begin{equation}\tag{\ref{eq:main-LR}}
\begin{gathered}
  \E_q\sup_{t\in\R}\norm{P_q[A(t),B]P_q}\le2\norm B\sum_{i=1}^k\sum_{y\in Y}F(|x_i-y|),\\
  \text{resp.}\quad\E_q\sup_{t\in\R}\norm{P_q[A(t),B]P_q}\le8|A|_1\norm B\sum_{x\in X,y\in Y}F(|x-y|).
\end{gathered}
\end{equation}
Consequently, let $\rho$ be a density matrix whose sector weights satisfy $\Tr(\rho P_q)\le K2^{-N}$ for all
$q$.  Then $\sup_t|\Tr(\rho[A(t),B])|$ is at most $K$ times the right side of \eqref{eq:main-LR}.  This covers $K=1$ and
$\rho=\ket{\Psi_0}\bra{\Psi_0}$, $\Psi_0=\ket{\uparrow\cdots\uparrow}_z\otimes\psi$, with \emph{arbitrary} $\psi$, as well
as $\rho=4^{-N}\one$.
\item \emph{(The operator norm sees the clean sector.)}  For all $N$, $x,y\in[1,N]$ and $t$,
\begin{equation}\tag{\ref{eq:main-LRnorm}}
\begin{gathered}
  \norm{\{\cdr_x(t),\cdr_y^\dagger\}}=\norm{[\cdr_x(t),n_y]}=\max_q|U_q(t)_{xy}|,\\
  \norm{[n_x(t),n_y]}=\max_q|U_q(t)_{xy}|\bigl(1-|U_q(t)_{xy}|^2\bigr)^{1/2}.
\end{gathered}
\end{equation}
The maximum includes the clean sector $q=\Zone=(1,\dots,1)$.  It has probability $2^{-N}$ under the sector
law, which vanishes as $N\to\infty$.  In particular, for $N\ge2$,
\[
  \sup_t\norm{\{\cdr_1(t),\cdr_N^\dagger\}}\ge\Bigl(\frac{3}{2(N+1)}\Bigr)^{1/2},\qquad
  \sup_t\norm{[n_1(t),n_N]}\ge\Bigl(\frac{3}{4(N+1)}\Bigr)^{1/2},
\]
whereas the sector average in \eqref{eq:main-prop} is at most $F(N-1)$.
\item \emph{(No operator-norm zero-velocity bound.)}  There is no function $G:\N_0\to[0,\infty)$ such that
$\sum_dd^2G(d)^2<\infty$ and $\sup_t\norm{[n_x(t),n_y]}\le G(|x-y|)$ for all $N\ge1$ and all $x,y\in[1,N]$.
The same holds with $[n_x(t),n_y]$ replaced by $\{\cdr_x(t),\cdr_y^\dagger\}$ or by $[\cdr_x(t),n_y]$.
\end{enumerate}
\end{stmtP}
\leanline{\texttt{avg\_\allowbreak{}norm\_\allowbreak{}evo\_\allowbreak{}le}, \texttt{avg\_\allowbreak{}iSup\_\allowbreak{}norm\_\allowbreak{}evo\_\allowbreak{}le}, \texttt{avg\_\allowbreak{}norm\_\allowbreak{}comm\_\allowbreak{}le} \,(\texttt{LiebRobinson.\allowbreak{}lean});\\ \texttt{LR\_\allowbreak{}sector\_\allowbreak{}mono}, \texttt{LR\_\allowbreak{}sector\_\allowbreak{}maj}, \texttt{LR\_\allowbreak{}avg\_\allowbreak{}mono}, \texttt{LR\_\allowbreak{}avg\_\allowbreak{}maj}, \texttt{LR\_\allowbreak{}trace\_\allowbreak{}mono}, \texttt{norm\_\allowbreak{}dressedAnticomm\_\allowbreak{}eq}, \texttt{norm\_\allowbreak{}comm\_\allowbreak{}number\_\allowbreak{}eq} \,(\texttt{LRFull.\allowbreak{}lean});\\ \texttt{numComm}, \texttt{norm\_\allowbreak{}comm\_\allowbreak{}nn\_\allowbreak{}eq}, \texttt{sup\_\allowbreak{}norm\_\allowbreak{}dressedAnticomm\_\allowbreak{}ge}, \texttt{sup\_\allowbreak{}norm\_\allowbreak{}comm\_\allowbreak{}nn\_\allowbreak{}ge}, \texttt{avg\_\allowbreak{}sup\_\allowbreak{}anticomm\_\allowbreak{}corner\_\allowbreak{}le} \,(\texttt{LRNorm.\allowbreak{}lean});\\ \texttt{no\_\allowbreak{}uniform\_\allowbreak{}LR\_\allowbreak{}nn}, \texttt{no\_\allowbreak{}uniform\_\allowbreak{}LR\_\allowbreak{}anticomm}, \texttt{no\_\allowbreak{}uniform\_\allowbreak{}LR\_\allowbreak{}cn} \,(\texttt{LRNoGo.\allowbreak{}lean})}
\renewcommand{\stmtname}{\Cref{thm:imbalance}}
\begin{stmtP}[Positive infinite-time imbalance for every $h\neq0$]
Let $N$ be even.
\begin{enumerate}[label=(\roman*)]
\item For every $h\in\R$, the limit $\Ibar:=\lim_{T\to\infty}T^{-1}\int_0^TI(t)\,dt$ exists and
$\Ibar=\E_q\,N^{-1}\sum_\alpha a_\alpha^2\in[0,1]$.
\item \emph{(Anticommutator identity.)}  For every $h\in\R$, $q$ and $\alpha$:
$E_\alpha a_\alpha=2h\,\langle\varphi_\alpha,V_q\Gamma\varphi_\alpha\rangle$.  If $h=0$, then $a_\alpha=0$ for all
$\alpha$ and $\Ibar=0$ (chiral obstruction).
\item \emph{(Unconditional bounds.)}  Let $h\neq0$.  For every block length $1\le L\le N$,
\begin{equation}\tag{\ref{eq:main-Ibound}}
  \Ibar\ \ge\ \frac{\lfloor N/L\rfloor}{N}\cdot\frac{h^2\bigl(\sqrt{L/2}-2\kappa_N\bigr)_+^2}{L\,(J+|h|)^2},
  \qquad\text{and}\qquad \Ibar\ \ge\ \frac{h^2}{N(J+|h|)^2}>0 .
\end{equation}
\item \emph{(Uniform in $N$.)}  Under Hypothesis~\ref{hyp:EC}, $\kappa_N\le\kappa:=2\SigF\le2M_1(F)$ for all
$N$, and for all even $N\ge2$
\begin{equation}\tag{\ref{eq:main-Iuniform}}
  \Ibar\ \ge\ \frac{h^2}{16\,L_*\,(J+|h|)^2},\qquad L_*:=\max\{2,\lceil32\kappa^2\rceil\}.
\end{equation}
\end{enumerate}
\end{stmtP}
\leanline{\texttt{stag\_\allowbreak{}hq\_\allowbreak{}anticomm}, \texttt{eigenvalue\_\allowbreak{}mul\_\allowbreak{}stagWeight}, \texttt{hq\_\allowbreak{}eigenvalues\_\allowbreak{}injective}, \texttt{imbOne\_\allowbreak{}timeAvg}, \texttt{avg\_\allowbreak{}imbInf\_\allowbreak{}ge}, \texttt{avg\_\allowbreak{}imbInf\_\allowbreak{}uniform}, \texttt{lgtImbalance\_\allowbreak{}eq}, \texttt{lgtImbalance\_\allowbreak{}timeAvg}, \texttt{lgtImbalance\_\allowbreak{}lower} \,(\texttt{Imbalance.\allowbreak{}lean}, \texttt{Main.\allowbreak{}lean})}
\renewcommand{\stmtname}{\Cref{cor:imbalanceopen}}
\begin{stmtP}[Imbalance on the open chain]
Consider the open chain of \cref{prop:variants}(i), i.e.\ the model of [15] without the boundary link,
with the density-wave matter state and $z$-polarized links, and let $\Ibar^{\rm open}$ be its infinite-time
imbalance.  Under Hypothesis~\ref{hyp:EC}, for all even $N\ge4$,
\[
  \Ibar^{\rm open}\ \ge\ \frac{h^2}{16\,L'_*\,(J+|h|)^2},\qquad L'_*:=\max\{2,\lceil128\kappa^2\rceil\},
\]
with $\kappa$ as in \cref{thm:imbalance}(iv).
\end{stmtP}
\leanline{\texttt{lgtImbalanceOpen\_\allowbreak{}lower}, \texttt{lgtImbalanceOpen\_\allowbreak{}timeAvg}, \texttt{lgtImbalanceOpen\_\allowbreak{}eq}, \texttt{avgPar\_\allowbreak{}imbInf\_\allowbreak{}uniform} \,(\texttt{OpenImbalance.\allowbreak{}lean})}
\renewcommand{\stmtname}{\Cref{thm:arealaw}}
\begin{stmtP}[Uniform area laws]
Let $\psi=\ket n$ be a Fock state and $1\le\ell<N$.
\begin{enumerate}[label=(\alph*)]
\item \emph{(Deterministic chain.)}  For every $q$ and $t$, with $T_q(t):=P_{[\ell+1,N]}K_q(t)P_{[1,\ell]}$,
\[
  S(\rho^q_A(t))\le2\ln2\,\norm{T_q(t)}_1\le2\ln2\sum_{x\le\ell<y}|K_q(t)_{xy}| .
\]
\item \emph{(Sector-averaged area law.)}  Under Hypothesis~\ref{hyp:EC},
\[
  \E_q\sup_{t\in\R}S(\rho^q_A(t))\le2\ln2\,\ThF\le8\ln2\,M_2(F),\qquad\text{hence}\qquad
  \sup_t\SG(t)\le8\ln2\,M_2(F),
\]
uniformly in $N$, $\ell$ and the Fock state $\psi$.
\item \emph{(Definite-charge initial states.)}  Let $\Psi_0=\ket{\xi}_\sigma\otimes\ket n$ with
$\xi\in\{\pm1\}^N$ (links in $\sigma^x$-eigenstates), and $q=q(n,\xi)$ as in \cref{lem:productlaw}(d).  Then
$\Sfull(t)=S(\rho^q_A(t))$ for all $t$.  If $\xi$ is uniformly distributed, so is $q$, and under
Hypothesis~\ref{hyp:EC}
\[
  \E_\xi\sup_t\Sfull(t)\le8\ln2\,M_2(F),\qquad
  \Prob_\xi\Bigl(\sup_t\Sfull(t)\ge s\Bigr)\le\frac{8\ln2\,M_2(F)}{s}\quad(s>0),
\]
and $\E_\xi\Sfull(t)=\SG(t)$ for every $t$.
\end{enumerate}
\end{stmtP}
\leanline{\texttt{vonNeumann\_\allowbreak{}slater\_\allowbreak{}le}, \texttt{lgt\_\allowbreak{}areaLaw} \,(\texttt{AreaLaw.\allowbreak{}lean});\\ \texttt{annealed\_\allowbreak{}K}, \texttt{annealed\_\allowbreak{}cut}, \texttt{vonNeumann\_\allowbreak{}slater\_\allowbreak{}sharp}, \texttt{lgt\_\allowbreak{}areaLaw\_\allowbreak{}sharp}, \texttt{thetaF\_\allowbreak{}le}, \texttt{lgt\_\allowbreak{}areaLaw\_\allowbreak{}sharp\_\allowbreak{}of\_\allowbreak{}ldt} \,(\texttt{AreaLawSharp.\allowbreak{}lean});\\ \texttt{vonNeumann\_\allowbreak{}rdmLGT\_\allowbreak{}xprod}, \texttt{lgt\_\allowbreak{}areaLaw\_\allowbreak{}full} \,(\texttt{CutLocal.\allowbreak{}lean});\\ \texttt{avg\_\allowbreak{}full\_\allowbreak{}eq\_\allowbreak{}SG}, \texttt{sup\_\allowbreak{}SG\_\allowbreak{}le}, \texttt{avg\_\allowbreak{}sup\_\allowbreak{}full\_\allowbreak{}le}, \texttt{prob\_\allowbreak{}sup\_\allowbreak{}full\_\allowbreak{}ge\_\allowbreak{}le} \,(\texttt{AreaLawFull.\allowbreak{}lean})}
\renewcommand{\stmtname}{\Cref{prop:entropies}}
\begin{stmtP}[Entropy taxonomy]
For every $t\in\R$:
\begin{enumerate}[label=(\alph*)]
\item \emph{(Duality defect.)}  $|\Sfull(t)-S_{A'}(\Phi(t))|\le\ln2$.
\item \emph{(Conditional versus mixed.)}  $\rho_{\rm mix}(t)$ is the marginal of $\rho_{A'}(\Phi(t))$ on the modes
$c_1,\dots,c_\ell$, and
\[
  \SG(t)\le\Smix(t)\le\min\{\SG(t)+H_t,\ \ell\ln2\},\qquad H_t\le N\ln2,\qquad H_0=0 .
\]
\item \emph{(Number entropy as a lower bound.)}  $p_t=\E_qp^q_t$, and for $X\in\{\mathrm{full},\mathrm{mix},f\}$
\[
  S_X(t)=\SN(t)+\sum_kp_t(k)\,S\bigl(\rho_X^{(k)}(t)\bigr)\ \ge\ \SN(t),
\]
where $p_t(k)\rho^{(k)}_X(t)$ is the block of the respective reduced state on $\{N_A=k\}$.  For the conditional
entropy one only has $\SG(t)\ge\E_q\Hsh(p^q_t)$, and $\E_q\Hsh(p^q_t)\le\SN(t)\le\E_q\Hsh(p^q_t)+H_t$.
\item \emph{(Superposition versus ensemble.)}  Let $\rho^{\rm ens}:=\sum_qP_q\ket{\Psi_0}\bra{\Psi_0}P_q$.  Then
\[
  \rho^{\rm ens}=2^{-N}\sum_{\xi\in\{\pm1\}^N}\ket{\xi}\bra{\xi}_\sigma\otimes\ket n\bra n=2^{-N}\one_{\rm links}\otimes\ket n\bra n,
\]
and for all gauge-invariant $O_1,\dots,O_m$ and $t_1,\dots,t_m\in\R$,
\begin{equation}\tag{\ref{eq:main-ens}}
  \langle\Psi_0,O_1(t_1)\cdots O_m(t_m)\Psi_0\rangle=\Tr\bigl(\rho^{\rm ens}O_1(t_1)\cdots O_m(t_m)\bigr).
\end{equation}
In particular, the following are the same for $\Psi(t)$ and for the ensemble
$\rho^{\rm ens}(t):=e^{-itH}\rho^{\rm ens}e^{itH}$: the law of every gauge-invariant observable; the restriction of
the state to the dressed-fermion algebra $\Alg([1,\ell])$, which is represented through $U$ by $\rho_{\rm mix}(t)$;
hence $\Smix(t)$; and $\SN(t)$.  The members $\ket\xi_\sigma\otimes\ket n$ of the
ensemble are definite-charge states, and their average full entanglement is $\SG(t)$
(\cref{thm:arealaw}(c)).
\end{enumerate}
\end{stmtP}
\leanline{\texttt{fullEntropy\_\allowbreak{}eq\_\allowbreak{}numberEntropy\_\allowbreak{}add}, \texttt{bareEntropy\_\allowbreak{}eq\_\allowbreak{}numberEntropy\_\allowbreak{}add}, \texttt{vonNeumann\_\allowbreak{}block\_\allowbreak{}decomp}, \texttt{lgt\_\allowbreak{}duality\_\allowbreak{}defect}, \texttt{mixRdm\_\allowbreak{}eq\_\allowbreak{}marginal}, \texttt{condEntropy\_\allowbreak{}le\_\allowbreak{}mix}, \texttt{mix\_\allowbreak{}le\_\allowbreak{}cond\_\allowbreak{}add}, \texttt{mix\_\allowbreak{}le\_\allowbreak{}log}, \texttt{sectorLawEntropy\_\allowbreak{}le}, \texttt{sectorLawEntropy\_\allowbreak{}zero}, \texttt{lgtNumDist\_\allowbreak{}eq\_\allowbreak{}avg}, \texttt{numberEntropy\_\allowbreak{}le\_\allowbreak{}mix}, \texttt{numberEntropy\_\allowbreak{}le\_\allowbreak{}bare}, \texttt{avg\_\allowbreak{}sector\_\allowbreak{}shannon\_\allowbreak{}le\_\allowbreak{}cond}, \texttt{avg\_\allowbreak{}sector\_\allowbreak{}shannon\_\allowbreak{}le\_\allowbreak{}numberEntropy}, \texttt{numberEntropy\_\allowbreak{}le\_\allowbreak{}avg\_\allowbreak{}add}, \texttt{shannon\_\allowbreak{}le\_\allowbreak{}vonNeumann\_\allowbreak{}of\_\allowbreak{}block}, \texttt{numDist\_\allowbreak{}shannon\_\allowbreak{}le\_\allowbreak{}vonNeumann\_\allowbreak{}slater}, \texttt{lgt\_\allowbreak{}numberEntropy\_\allowbreak{}le\_\allowbreak{}full} \,(\texttt{Entropies.\allowbreak{}lean}, \texttt{Taxonomy.\allowbreak{}lean}, \texttt{EntropyBlocks.\allowbreak{}lean});\\ \texttt{expect\_\allowbreak{}initZ\_\allowbreak{}eq\_\allowbreak{}trace\_\allowbreak{}ensemble}, \texttt{correlator\_\allowbreak{}eq\_\allowbreak{}ensemble}, \texttt{gaugeInv\_\allowbreak{}heis} \,(\texttt{Ensemble.\allowbreak{}lean})}
\renewcommand{\stmtname}{\Cref{lem:intentropy}}
\begin{stmtP}[Entropy of integer variables {[40]}]
Let $X$ be an integer-valued random variable taking finitely many values.  Then
$\Hsh(X)\le\frac12\ln\bigl(2\pi e(\Var X+\frac1{12})\bigr)$.
\end{stmtP}
\leanline{\texttt{massey}, \texttt{shannon\_\allowbreak{}le\_\allowbreak{}log\_\allowbreak{}card}, \texttt{pvar\_\allowbreak{}mixture}, \texttt{binEntropy\_\allowbreak{}le\_\allowbreak{}two\_\allowbreak{}sqrt} \,(\texttt{Entropy.\allowbreak{}lean});\\ \texttt{wStateS}, \texttt{energyDensityS}, \texttt{lawEntropy}, \texttt{wilson\_\allowbreak{}field\_\allowbreak{}bound\_\allowbreak{}super}, \texttt{bgEntropyS\_\allowbreak{}wStateS}, \texttt{wilson\_\allowbreak{}numberEntropy\_\allowbreak{}le\_\allowbreak{}super}, \texttt{lawEntropy\_\allowbreak{}le\_\allowbreak{}massey}, \texttt{lawEntropy\_\allowbreak{}le\_\allowbreak{}log\_\allowbreak{}card}, \texttt{wilson\_\allowbreak{}numberEntropy\_\allowbreak{}avg\_\allowbreak{}super}, \texttt{bgEntropyS\_\allowbreak{}bare\_\allowbreak{}vacuum}, \texttt{bgEntropyS\_\allowbreak{}le\_\allowbreak{}log\_\allowbreak{}three} \,(\texttt{WilsonSuper.\allowbreak{}lean})}
\renewcommand{\stmtname}{\Cref{thm:numberentropy}}
\begin{stmtP}[Bounded number entropy]
Assume Hypothesis~\ref{hyp:EC}.  Let $\psi$ be a Fock state, $\Psi_0=\ket{\uparrow\cdots\uparrow}_z\otimes\psi$ and
$A=[1,\ell]$, $1\le\ell<N$.  Then, for all $t\in\R$ and all $N$,
\begin{equation}\tag{\ref{eq:main-SN}}
  \SN(t)\le\tfrac12\ln\bigl(2\pi e\,(K_1+K_2+\tfrac1{12})\bigr),\qquad
  K_1:=\ThF\le4M_2(F),\quad K_2:=\SigpF^2\le M'_1(F)^2 .
\end{equation}
Moreover $\E_q\sup_t\Hsh(p^q_t)\le\frac12\ln\bigl(2\pi e(K_1+\frac1{12})\bigr)$ for the number entropy in a single sector.
\end{stmtP}
\leanline{\texttt{dressed\_\allowbreak{}anticomm\_\allowbreak{}t}, \texttt{norm\_\allowbreak{}sector\_\allowbreak{}anticomm}, \texttt{avg\_\allowbreak{}sup\_\allowbreak{}anticomm\_\allowbreak{}le}, \texttt{expect\_\allowbreak{}abs\_\allowbreak{}anticomm\_\allowbreak{}le}, \texttt{one\_\allowbreak{}particle\_\allowbreak{}moments}, \texttt{sector\_\allowbreak{}transport\_\allowbreak{}le}, \texttt{avg\_\allowbreak{}sup\_\allowbreak{}transport\_\allowbreak{}le}, \texttt{lgt\_\allowbreak{}transport\_\allowbreak{}sharp}, \texttt{lgt\_\allowbreak{}transport\_\allowbreak{}sq}, \texttt{dressed\_\allowbreak{}corr\_\allowbreak{}le}, \texttt{lgt\_\allowbreak{}numberEntropy\_\allowbreak{}sharp}, \texttt{avg\_\allowbreak{}sup\_\allowbreak{}sector\_\allowbreak{}numberEntropy\_\allowbreak{}le} \,(\texttt{TransportSharp.\allowbreak{}lean});\\ \texttt{lgt\_\allowbreak{}numberEntropy}, \texttt{lgtNumDist\_\allowbreak{}isProb} \,(\texttt{Main.\allowbreak{}lean})}
\renewcommand{\stmtname}{\Cref{thm:infvol}}
\begin{stmtP}[Thermodynamic limit; unconditional]
Let $\lambda\neq0$.  In the setting above:
\begin{enumerate}[label=(\alph*)]
\item \emph{(Limit dynamics.)}  For every finite interval $X$, $O\in\GAlg(X)$ and $t\in\R$, the limit
\[
  \omega_t(O):=\lim_{M\to\infty}\langle\Psi^{(M)}(t),O\,\Psi^{(M)}(t)\rangle=\E\,\Tr_{\cF_X}\bigl(\rhoqf{K^\Z_q(t)|_X}\,O_q\bigr)
\]
exists, where $K|_X:=P_XKP_X$ on $\ell^2(X)$.  The convergence holds sector by sector, uniformly in $q$, and
uniformly for $t$ in compact sets.
\item \emph{(Propagator, moments, transport.)}  For all $x,y\in\Z$, $\E\sup_t|U^\Z_q(t)_{xy}|\le\FZ(|x-y|)$, and
$\E|U^{(M)}_q(t)_{xy}|\to\E|U^\Z_q(t)_{xy}|$ for every $t$.  If $S=\{y\}$ (one particle), then
$\langle\Psi^{(M)}(t),n_x\Psi^{(M)}(t)\rangle\to\E|U^\Z_q(t)_{xy}|^2$, and for $0\le p\le2$
\[
  \sup_t\sum_{x\in\Z}(|x-y|+1)^p\,\E|U^\Z_q(t)_{xy}|^2\le2M_p(\FZ).
\]  For general $S$ and every $m\in\Z$, the
transported charge
\[
  \mathcal T_m(t):=\lim_{M\to\infty}\Bigl(\langle\Psi^{(M)}(t),N_{[-M,m]}\Psi^{(M)}(t)\rangle-|S\cap[-M,m]|\Bigr)
\]
exists for every $t$, and $\sup_t|\mathcal T_m(t)|\le\Sigma_1(\FZ)$.
\item \emph{(Lieb--Robinson.)}  Let $X\subset\Z$ be finite and $Y\subset\Z$ a finite interval.  Let $\Alg^+(X)$ be the
even part of the algebra generated by the $\cdr^{(M)}_x$, $x\in X$; it is spanned by products of the even monomials
$\cdr^{(M)\sharp}_x\cdr^{(M)\sharp}_y$, which do not depend on $M$.  Then for $A\in\Alg^+(X)$, $B\in\GAlg(Y)$ and every $t$,
\[
  \limsup_{M\to\infty}\E\norm{P_q[A(t),B]}\le8|A|_1\norm B\sum_{x\in X,y\in Y}\FZ(|x-y|),
\]
and likewise with the monomial bound of \eqref{eq:main-LR}.  The right side does not depend on $t$.
\item \emph{(Area law for bulk intervals.)}  For a finite interval $A\subset\Z$, let $S^{(M)}_q(t)$ be the entropy of
the restriction of $\psi^{(M)}_q(t)$ to the $c$-modes in $A$ (\cref{lem:slater}), $\SG^{(M)}(t):=\E S^{(M)}_q(t)$ the
bulk conditional entanglement, and $S^\Z_q(t):=S(\rhoqf{K^\Z_q(t)|_A})$.  For every $t$, $S^{(M)}_q(t)\to S^\Z_q(t)$
uniformly in $q$, so $\SG^{(M)}(t)\to\SG^\Z(t):=\E S^\Z_q(t)$.  Moreover
\[
  \E\sup_tS^\Z_q(t)\le4\ln2\,\Theta(\FZ)\le16\ln2\,M_2(\FZ).
\]
For the half-infinite regions $A_M$ (sites $-M,\dots,m$ and links $b_{-M},\dots,b_m$) one has, for every $t$,
$\limsup_M\E S(\rho^{q}_{A_M}(t))\le2\ln2\,\Theta(\FZ)$.  The same bound holds for $\limsup_M\E_\xi\Sfull^{(M)}(t)$,
for definite-charge initial states $\ket{\xi}_\sigma\otimes\ket{n|_{\Lambda_M}}$ with $\xi$ uniform.
\item \emph{(Number entropy for bulk intervals.)}  For a finite interval $A\subset\Z$ and every $t$, the law of $N_A$
in $\Psi^{(M)}(t)$ converges to $p_t:=\E\,\mathrm{PB}(\nu(K^\Z_q(t)|_A))$, where $\mathrm{PB}(\nu)$ denotes the
Poisson-binomial law with the eigenvalues $\nu$ of $K^\Z_q(t)|_A$.  Hence $\SN^{(M)}(t)\to\SN(t):=\Hsh(p_t)$, and
\[
  \sup_t\SN(t)\le\tfrac12\ln\bigl(2\pi e\,(2\Theta(\FZ)+4\Sigma_1'(\FZ)^2+\tfrac1{12})\bigr).
\]
\end{enumerate}
In (b)--(e) the bounds concern the limit objects, for every $t$; we do not claim that the limits $M\to\infty$ and
$\sup_t$ can be interchanged.  (By Theorem~\ref{thm:EC}, which is not used here, the finite-volume bounds also hold
uniformly in $M$.)
\end{stmtP}
\leanline{\texttt{rdmF}, \texttt{reorderSign}, \texttt{vonNeumann\_\allowbreak{}rdmF\_\allowbreak{}slater}, \texttt{vonNeumann\_\allowbreak{}rdmF\_\allowbreak{}le}, \texttt{vonNeumann\_\allowbreak{}rdmF\_\allowbreak{}segment}, \texttt{lgt\_\allowbreak{}areaLaw\_\allowbreak{}bulk} \,(\texttt{SlaterBulk.\allowbreak{}lean});\\ \texttt{lgt\_\allowbreak{}bulk\_\allowbreak{}transport\_\allowbreak{}sq}, \texttt{lgt\_\allowbreak{}numberEntropy\_\allowbreak{}bulk} \,(\texttt{NumberEntropyBulk.\allowbreak{}lean});\\ \texttt{tendsto\_\allowbreak{}evoZ}, \texttt{lintegral\_\allowbreak{}iSup\_\allowbreak{}UZ\_\allowbreak{}le}, \texttt{infvol\_\allowbreak{}propagator\_\allowbreak{}of\_\allowbreak{}ldt}, \texttt{tendsto\_\allowbreak{}density\_\allowbreak{}one}, \texttt{moments\_\allowbreak{}UZ}, \texttt{abs\_\allowbreak{}transportM\_\allowbreak{}le}, \texttt{transport\_\allowbreak{}tendsto} \,(\texttt{InfVol.\allowbreak{}lean}, \texttt{InfVolTransport.\allowbreak{}lean});\\ \texttt{KZ}, \texttt{KZA}, \texttt{pZ}, \texttt{tendsto\_\allowbreak{}numDistM}, \texttt{tendsto\_\allowbreak{}numberEntropyM}, \texttt{numberEntropyZ\_\allowbreak{}le} \,(\texttt{InfVolNumber.\allowbreak{}lean});\\ \texttt{majListM}, \texttt{infvol\_\allowbreak{}LR} \,(\texttt{InfVolLR.\allowbreak{}lean});\\ \texttt{embF}, \texttt{expect\_\allowbreak{}embF}, \texttt{locObsM}, \texttt{KZL}, \texttt{tendsto\_\allowbreak{}omega}, \texttt{tendstoUniformlyOn\_\allowbreak{}sector} \,(\texttt{InfVolDyn.\allowbreak{}lean});\\ \texttt{entM}, \texttt{entZ}, \texttt{tendstoUniformly\_\allowbreak{}entM}, \texttt{tendsto\_\allowbreak{}SGM}, \texttt{lintegral\_\allowbreak{}iSup\_\allowbreak{}entZ\_\allowbreak{}le}, \texttt{areaLaw\_\allowbreak{}halfInfinite} \,(\texttt{InfVolArea.\allowbreak{}lean})}
\renewcommand{\stmtname}{\Cref{lem:duhamel}}
\begin{stmtP}[Duhamel sensitivity to distant charges]
Let $q,q'\in\{\pm1\}^N$ differ exactly on $D\subset[1,N]$, and let $z\in[1,N]$, $s\in\R$.  Then
\[
  \norm{(U_q(s)-U_{q'}(s))\delta_z}\le4|h|\int_0^{|s|}\norm{\one_DU_{q'}(\pm r)\delta_z}\,dr\le4|h|\,|s|\sum_{x\in D}Q_{q'}(x,z),
\]
with the sign of $s$ in $\pm r$.  Under Hypothesis~\ref{hyp:EC}, with $q^D$ the configuration $q$ flipped on $D$,
\[
  \E_q\sup_{|s|\le t}\norm{(U_q(s)-U_{q^D}(s))\delta_z}\le4|h|\,t\sum_{x\in D}F(|x-z|).
\]
In particular, if $\dist(z,D)\ge R$ and $F(d)=\min\{1,Ce^{-\eta d}\}$, the right side is at most
$8|h|tCe^{-\eta R}/(1-e^{-\eta})$.
\end{stmtP}
\leanline{\texttt{colNorm\_\allowbreak{}evo\_\allowbreak{}sub\_\allowbreak{}le}, \texttt{avg\_\allowbreak{}colNorm\_\allowbreak{}evo\_\allowbreak{}sub\_\allowbreak{}le} \,(\texttt{Duhamel.\allowbreak{}lean})}
\renewcommand{\stmtname}{\Cref{conj:logent}}
\begin{stmtP}[Logarithmic growth of the full entanglement]
Let $\psi$ be the CDW, $\Psi_0=\ket{\uparrow\cdots\uparrow}_z\otimes\psi$ and $\ell=\lfloor N/2\rfloor$.  There are constants
$C_1,C_2<\infty$, depending only on $J$ and $h$, such that $\Sfull(t)\le C_1\ln(2+|t|)+C_2$ for all $N$ and $t$.  Moreover
$\lim_{t\to\infty}\liminf_{N\to\infty}\Sfull(t)=\infty$.
\end{stmtP}
\leanline{conjecture, not proved (not formalized).}
\renewcommand{\stmtname}{\Cref{conj:volmix}}
\begin{stmtP}[Volume law for the mixed entropies]
There is $c=c(J,h)>0$ such that, for all $N\ge2\ell$ and the CDW initial state,
$\liminf_{T\to\infty}T^{-1}\int_0^T\Smix(t)\,dt\ge c\,\ell$, and likewise for $\Sfm$.
\end{stmtP}
\leanline{conjecture, not proved (not formalized).}
\subsection*{\S\,\ref{sec:dimer}\quad Disorder-free localization depends on the initial state}
\addcontentsline{toc}{subsection}{\S\,\ref{sec:dimer}\quad Disorder-free localization depends on the initial state}
\renewcommand{\stmtname}{\Cref{def:dimerstate}}
\begin{stmtD}[Dimer product states]
Let $N=2M$.  For $n\in\{0,1\}^N$ put $\xi_{2k}(n):=(-1)^{N_{\le2k}(n)}$, $k=1,\dots,M$, and
\begin{equation}\tag{\ref{eq:dimerstate1}}
  \chidim(n):=\bigotimes_{k=1}^M\Bigl(\ket{\uparrow}_{b_{2k-1}}\otimes\ket{\xi_{2k}(n)}_{b_{2k}}\Bigr)\in\LinkS_N,
  \qquad \Psidim(n):=\ket{n}\otimes\chidim(n)\in\cH_N,
\end{equation}
where $\ket{\pm1}_{b}$ denotes the $\sigma^x_b$-eigenvector $\ket{\pm}$ of \eqref{eq:szflip}.  The odd
links are $z$-polarized, as in [15]; the even links are $x$-polarized, with orientation
given by the matter parity to their left.  $\Omdim:=\Psidim(0)$ (no fermions, all even links in
$\ket{+}$) is the \emph{dimer vacuum}, and $\Psi^x:=\Psidim(e_x)$, $x\in\{1,\dots,N\}$, is the state
with a single fermion at $x$.  For the charge-density wave $\ncdw$ ($n_x=1$ iff $x$ is even) one has
$\xi_{2k}(\ncdw)=(-1)^k$, and $\Psidim(\ncdw)$ is the state of \cref{lem:productlaw}(e).
\end{stmtD}
\leanline{\texttt{dimerState}, \texttt{prodStateXi\_\allowbreak{}dual}, \texttt{pairDefect\_\allowbreak{}eq\_\allowbreak{}zero\_\allowbreak{}iff}, \texttt{gauss\_\allowbreak{}pair\_\allowbreak{}prodStateXi}, \texttt{dimerState\_\allowbreak{}dual}, \texttt{normSq\_\allowbreak{}dimerState}, \texttt{normSq\_\allowbreak{}gaussProj\_\allowbreak{}dimerState}, \texttt{dimerState\_\allowbreak{}cdw}, \texttt{gauss\_\allowbreak{}pair\_\allowbreak{}dimerState}, \texttt{gauss\_\allowbreak{}prefix\_\allowbreak{}dimerState}, \texttt{dressed\_\allowbreak{}cre\_\allowbreak{}dimerState}, \texttt{dimerState\_\allowbreak{}eq\_\allowbreak{}dressed}, \texttt{dimerState\_\allowbreak{}single\_\allowbreak{}eq}, \texttt{sector\_\allowbreak{}average\_\allowbreak{}dimerState}, \texttt{dimer\_\allowbreak{}density\_\allowbreak{}single}, \texttt{dimerInit\_\allowbreak{}dual}, \texttt{card\_\allowbreak{}isDimer}, \texttt{sector\_\allowbreak{}average\_\allowbreak{}dimer} \,(\texttt{Dimer.\allowbreak{}lean}, \texttt{DimerState.\allowbreak{}lean})}
\renewcommand{\stmtname}{\Cref{prop:dimerlaw}}
\begin{stmtP}[Dimer sector law]
Let $N=2M$ and $n\in\{0,1\}^N$.  For $\eta\in\{\pm1\}^M$ let
$q(\eta):=(\eta_1,\eta_1,\eta_2,\eta_2,\dots,\eta_M,\eta_M)\in\{\pm1\}^N$.
\begin{enumerate}[label=(\alph*)]
\item $U^\dagger\Psidim(n)=2^{-M/2}\sum_{\eta\in\{\pm1\}^M}\ket{n}\otimes\ket{q(\eta)}$.  Hence
$\norm{P_q\Psidim(n)}^2=2^{-M}$ if $q=q(\eta)$ for some $\eta$, and $\norm{P_q\Psidim(n)}^2=0$
otherwise.  The sector law is the \emph{random dimer law} $\mudim$: $(q_{2k-1},q_{2k})=(\eta_k,\eta_k)$
with $\eta_1,\dots,\eta_M$ i.i.d.\ uniform on $\{\pm1\}$.  Every sector in its support carries the
matter state $\ket{n}$, with the same positive amplitude.  (For $n=\ncdw$ this is \eqref{eq:dimerstate}.)
\item \emph{(Neutral dimers.)}  $\prod_{j\le2k}G_j\,\Psidim(n)=\Psidim(n)$ and
$G_{2k-1}G_{2k}\,\Psidim(n)=\Psidim(n)$ for $k=1,\dots,M$.
\item \emph{(Dressed fermions.)}  Let $s_1<\dots<s_m$ be the occupied sites of $n$.  Then
\[
  \Psidim(n)=\cdr_{s_1}^\dagger\cdots\cdr_{s_m}^\dagger\,\Omdim;\qquad\text{in particular}\quad\Psi^x=\cdr_x^\dagger\Omdim .
\]
\item \emph{(Dynamics.)}  Let $\Psi(t):=\e^{-\ii tH}\Psidim(n)$ and let $\Edim$ denote the expectation
with respect to $\mudim$.  For every gauge-invariant $O$, with $O_q$ as in \eqref{eq:Oq},
\begin{equation}\tag{\ref{eq:dimeravg}}
  \langle\Psi(t),O\,\Psi(t)\rangle=\Edim\bigl\langle\Gamma(\e^{-\ii t\hq})\,n,\;O_q\,\Gamma(\e^{-\ii t\hq})\,n\bigr\rangle .
\end{equation}
In particular, for $\Psi^x(t):=\e^{-\ii tH}\Psi^x$ and all $y$,
\begin{equation}\tag{\ref{eq:dimerdensity}}
  \langle\Psi^x(t),n_y\,\Psi^x(t)\rangle=\Edim\bigl|\langle\delta_y,\e^{-\ii t\hq}\delta_x\rangle\bigr|^2 .
\end{equation}
\item \emph{(Other orientations of the even links.)}  Let $\xi=(\xi_2,\xi_4,\dots,\xi_{2M})\in\{\pm1\}^M$
be arbitrary, $\xi_0:=1$, and
$\chi_\xi:=\bigotimes_{k=1}^M(\ket{\uparrow}_{b_{2k-1}}\otimes\ket{\xi_{2k}}_{b_{2k}})$.  Then
$U^\dagger(\ket{n}\otimes\chi_\xi)=2^{-M/2}\sum_{\epsilon\in\{\pm1\}^M}\ket{n}\otimes\ket{q^\epsilon}$ with
\begin{equation}\tag{\ref{eq:dimersigns}}
  q^\epsilon_{2k-1}=s_k\epsilon_k,\qquad q^\epsilon_{2k}=d_k\,s_k\epsilon_k,\qquad
  s_k:=\xi_{2k-2}(-1)^{n_{2k-1}},\qquad d_k:=(-1)^{n_{2k-1}+n_{2k}}\,\xi_{2k-2}\,\xi_{2k} .
\end{equation}
The sector law is the law of $(\eta_k,d_k\eta_k)_{k\le M}$ with $\eta$ i.i.d.\ uniform.  It equals
$\mudim$ if and only if $d_k=1$ for all $k$, i.e.\ if and only if $\xi_{2k}=\xi_{2k}(n)$ for all $k$.
On $\ket{n}\otimes\chi_\xi$ the operator $G_{2k-1}G_{2k}$ acts as the number $d_k$.
\end{enumerate}
\end{stmtP}
\leanline{\texttt{dimerState}, \texttt{prodStateXi\_\allowbreak{}dual}, \texttt{pairDefect\_\allowbreak{}eq\_\allowbreak{}zero\_\allowbreak{}iff}, \texttt{gauss\_\allowbreak{}pair\_\allowbreak{}prodStateXi}, \texttt{dimerState\_\allowbreak{}dual}, \texttt{normSq\_\allowbreak{}dimerState}, \texttt{normSq\_\allowbreak{}gaussProj\_\allowbreak{}dimerState}, \texttt{dimerState\_\allowbreak{}cdw}, \texttt{gauss\_\allowbreak{}pair\_\allowbreak{}dimerState}, \texttt{gauss\_\allowbreak{}prefix\_\allowbreak{}dimerState}, \texttt{dressed\_\allowbreak{}cre\_\allowbreak{}dimerState}, \texttt{dimerState\_\allowbreak{}eq\_\allowbreak{}dressed}, \texttt{dimerState\_\allowbreak{}single\_\allowbreak{}eq}, \texttt{sector\_\allowbreak{}average\_\allowbreak{}dimerState}, \texttt{dimer\_\allowbreak{}density\_\allowbreak{}single}, \texttt{dimerInit\_\allowbreak{}dual}, \texttt{card\_\allowbreak{}isDimer}, \texttt{sector\_\allowbreak{}average\_\allowbreak{}dimer} \,(\texttt{Dimer.\allowbreak{}lean}, \texttt{DimerState.\allowbreak{}lean})}
\renewcommand{\stmtname}{\Cref{lem:dimercritical}}
\begin{stmtP}[Critical energies of the dimer model]
Let $\lambda>0$.
\begin{enumerate}[label=(\roman*)]
\item For $E\in\R$, $[D_+(E),D_-(E)]=2\lambda(\lambda^2-E^2)\,\sigma_x$.  So $D_+(E)$ and $D_-(E)$ commute if
and only if $E=\pm\lambda$.
\item $E\in\R$ is a critical energy in the sense of [29, Def.~1] if and only if
$E\in\{\lambda,-\lambda\}$ and $\lambda<1$.  In particular, the dimer model has no critical energy when
$\lambda\ge1$, i.e.\ when $2h\ge J$.
\item Let $0<\lambda<1$, let $\theta\in(\pi/2,\pi)$ be defined by $\cos\theta=-\lambda$, and let
$S:=\bigl(\begin{smallmatrix}\sin\theta&\cos\theta\\0&1\end{smallmatrix}\bigr)$.  Then
\begin{equation}\tag{\ref{eq:dimerrot}}
  S^{-1}D_+(\lambda)S=-\one=\Rot_\pi,\qquad S^{-1}\Tdim_{-2\lambda}S=\Rot_\theta,\qquad S^{-1}D_-(\lambda)S=\Rot_{2\theta}.
\end{equation}
Thus $S^{-1}$ is a matrix $M$ as in [29, Eq.~(7)] for the critical energy $E_c=\lambda$, with
rotation angles $\eta_+=\pi$ and $\eta_-=2\theta$.  With $\langle c_\pm\rangle:=p_+c_++p_-c_-$,
\begin{equation}\tag{\ref{eq:dimerangles}}
  \langle\e^{2\ii\eta_\pm}\rangle=\tfrac12\bigl(1+\e^{4\ii\theta}\bigr),\qquad
  \bigl|\langle\e^{2\ii\eta_\pm}\rangle\bigr|=|1-2\lambda^2|<1,\qquad
  \langle\e^{4\ii\eta_\pm}\rangle=1\iff\lambda=1/\sqrt2 .
\end{equation}
Moreover $\kappa:=\norm{S}\,\norm{S^{-1}}=\bigl((1+\lambda)/(1-\lambda)\bigr)^{1/2}$.
\item The same holds at $E_c=-\lambda$ with the roles of $D_+$ and $D_-$ exchanged: with
$\Gamma_2:=\diag(1,-1)$, $(\Gamma_2S)^{-1}D_-(-\lambda)\,\Gamma_2S=\Rot_\pi$ and
$(\Gamma_2S)^{-1}D_+(-\lambda)\,\Gamma_2S=\Rot_{2\theta}$, so that $\eta_+=2\theta$, $\eta_-=\pi$, and
\eqref{eq:dimerangles} holds unchanged.
\end{enumerate}
\end{stmtP}
\leanline{\texttt{dimerTM\_\allowbreak{}plus\_\allowbreak{}crit}, \texttt{dimerTM\_\allowbreak{}minus\_\allowbreak{}elliptic}, \texttt{dimerTM\_\allowbreak{}commute\_\allowbreak{}crit}, \texttt{jsbs\_\allowbreak{}condition}, \texttt{dimerTM\_\allowbreak{}minus\_\allowbreak{}not\_\allowbreak{}elliptic} \,(\texttt{DimerCritical.\allowbreak{}lean});\\ \texttt{dimerD\_\allowbreak{}commutator}, \texttt{dimerD\_\allowbreak{}commute\_\allowbreak{}iff}, \texttt{dimerTheta\_\allowbreak{}spec}, \texttt{dimerS\_\allowbreak{}conj\_\allowbreak{}plus}, \texttt{dimerS\_\allowbreak{}conj\_\allowbreak{}tm}, \texttt{dimerS\_\allowbreak{}conj\_\allowbreak{}minus}, \texttt{dimer\_\allowbreak{}angles}, \texttt{dimerS\_\allowbreak{}kappa}, \texttt{gamma2\_\allowbreak{}tmC}, \texttt{dimerS\_\allowbreak{}conj\_\allowbreak{}neg}, \texttt{dimer\_\allowbreak{}transfer\_\allowbreak{}bound} \,(\texttt{DimerTransfer.\allowbreak{}lean})}
\renewcommand{\stmtname}{\Cref{lem:dimertransfer}}
\begin{stmtP}[Uniform transfer-matrix bound near $E_c=\lambda$]
Let $0<\lambda<1$, let $\kappa$ be as in \cref{lem:dimercritical}(iii), and put
\[
  K_1:=\kappa\,(2+2\lambda)^2,\qquad K_2:=\kappa\,(3+4\lambda)/2 .
\]
For every dimer configuration $v$ on an interval $I$, of either alignment, every $z\in\C$ with
$|z-\lambda|\le1$, and all $m\le n$ with $[m,n-1]\subset I$,
\begin{equation}\tag{\ref{eq:dimerTbound}}
  \norm{\Tdim^z(n,m)}\le K_1\,\e^{K_2|z-\lambda|\,(n-m)} .
\end{equation}
\end{stmtP}
\leanline{\texttt{dimerD\_\allowbreak{}commutator}, \texttt{dimerD\_\allowbreak{}commute\_\allowbreak{}iff}, \texttt{dimerTheta\_\allowbreak{}spec}, \texttt{dimerS\_\allowbreak{}conj\_\allowbreak{}plus}, \texttt{dimerS\_\allowbreak{}conj\_\allowbreak{}tm}, \texttt{dimerS\_\allowbreak{}conj\_\allowbreak{}minus}, \texttt{dimer\_\allowbreak{}angles}, \texttt{dimerS\_\allowbreak{}kappa}, \texttt{gamma2\_\allowbreak{}tmC}, \texttt{dimerS\_\allowbreak{}conj\_\allowbreak{}neg}, \texttt{dimer\_\allowbreak{}transfer\_\allowbreak{}bound} \,(\texttt{DimerTransfer.\allowbreak{}lean})}
\renewcommand{\stmtname}{\Cref{prop:dimerlower}}
\begin{stmtP}[Every dimer configuration spreads at least diffusively]
Let $0<\lambda<1$, $p>0$, let $K_1,K_2$ be as in \cref{lem:dimertransfer}, and
\begin{equation}\tag{\ref{eq:clp}}
  c_{\lambda,p}:=\frac{c_1^2}{2^{p+3}\,\pi\,(p+1)},\qquad c_1:=\frac{1}{(3+2\lambda)\,K_1\,\e^{K_2}} .
\end{equation}
Let $v$ be a dimer configuration on an interval $I\subset\Z$, let $x\in I$, and let $T\ge1$ with
$[x-T,x+T]\subset I$.  Then
\begin{equation}\tag{\ref{eq:dimerlower}}
  \Mom^I_{v,p}(T;x)\ge c_{\lambda,p}\,T^{p-1}.
\end{equation}
In particular, \eqref{eq:dimerlower} holds for every dimer configuration on $\Z$, every $x\in\Z$ and
every $T\ge1$.
\end{stmtP}
\leanline{\texttt{dimerMom}, \texttt{dimerGreen}, \texttt{dimerMom\_\allowbreak{}plancherel}, \texttt{dimer\_\allowbreak{}lower\_\allowbreak{}bound}, \texttt{dimerMomZ}, \texttt{tendsto\_\allowbreak{}dimerMom\_\allowbreak{}Z}, \texttt{dimer\_\allowbreak{}lower\_\allowbreak{}bound\_\allowbreak{}Z} \,(\texttt{DimerLower.\allowbreak{}lean}, \texttt{DimerInfVol.\allowbreak{}lean})}
\renewcommand{\stmtname}{\Cref{lem:dimerTL}}
\begin{stmtP}[Finite speed and the thermodynamic limit]
Let $\Lambda\subset\Lambda'\subset\Z$ be intervals containing $0$ ($\Lambda'=\Z$ allowed), let
$w':\Lambda'\to[-W,W]$, $w:=w'|_\Lambda$, $H_\Lambda:=-\Delta_\Lambda+\diag(w)$,
$H_{\Lambda'}:=-\Delta_{\Lambda'}+\diag(w')$, and $R:=\dist(0,\Z\setminus\Lambda)\in\N\cup\{\infty\}$.  For
all $s\in\R$ and $y\in\Lambda$:
\begin{enumerate}[label=(\roman*)]
\item $|\langle\delta_y,\e^{-\ii sH_\Lambda}\delta_0\rangle|\le\sum_{k\ge|y|}\frac{(2|s|)^k}{k!}\le\frac{(2|s|)^{|y|}}{|y|!}\,\e^{2|s|}$;
\item $|\langle\delta_y,(\e^{-\ii sH_\Lambda}-\e^{-\ii sH_{\Lambda'}})\delta_0\rangle|\le2\sum_{k\ge R}\frac{((2+W)|s|)^k}{k!}$;
\item if $\Lambda$ is finite and $0<p\le2$, then $\sum_{y\in\Lambda}|y|^p|\langle\delta_y,\e^{-\ii sH_\Lambda}\delta_0\rangle|^2\le(2|s|)^p$.
\end{enumerate}
\end{stmtP}
\leanline{\texttt{dimerTL\_\allowbreak{}i}, \texttt{dimerTL\_\allowbreak{}i{\textquotesingle}}, \texttt{dimerTL\_\allowbreak{}ii}, \texttt{dimerTL\_\allowbreak{}ii\_\allowbreak{}Z}, \texttt{dimerTL\_\allowbreak{}iii}, \texttt{evoZw}, \texttt{moment\_\allowbreak{}poly\_\allowbreak{}bound}, \texttt{tendsto\_\allowbreak{}moment\_\allowbreak{}Z} \,(\texttt{DimerTL.\allowbreak{}lean}, \texttt{DimerInfVol.\allowbreak{}lean})}
\renewcommand{\stmtname}{\Cref{lem:dimersecond}}
\begin{stmtP}[Second moment near the critical energy]
Let $0<\lambda<1$.  There are $C,K,\delta_0>0$, depending only on $\lambda$, such that for every $z\in\C$ with
$|z-\lambda|\le\delta_0$ and every $n\in\N$,
\[
  2^{-n}\sum_{\epsilon\in\{+,-\}^n}\norm{D_{\epsilon_n}(z)\cdots D_{\epsilon_1}(z)}_2^2\le C\,\e^{Kn((\Re z-\lambda)^2+|\Im z|)},
\]
where $\norm\cdot_2$ is the Hilbert--Schmidt norm.
\end{stmtP}
\leanline{\texttt{hs2}, \texttt{dimer\_\allowbreak{}second\_\allowbreak{}moment}, \texttt{integral\_\allowbreak{}hs2\_\allowbreak{}tmProd\_\allowbreak{}le}, \texttt{dimer\_\allowbreak{}super\_\allowbreak{}annealed}, \texttt{dimer\_\allowbreak{}super\_\allowbreak{}ae} \,(\texttt{DimerSecond.\allowbreak{}lean}, \texttt{DimerSuper.\allowbreak{}lean})}
\renewcommand{\stmtname}{\Cref{prop:dimersuper}}
\begin{stmtP}[Superdiffusion]
Let $0<\lambda<1$, $0<p\le2$ and $r\in\{0,1\}$, and let $v:=v^{(r)}$.  There are $c,T_0>0$, depending only on $\lambda$
and $p$, such that:
\begin{enumerate}[label=(\alph*)]
\item $\E\,\Mom^\Z_{v,p}(T;0)\ge c\,T^{p-1/2}$ for all $T\ge T_0$;
\item almost surely, for every $\alpha>0$, $\Mom^\Z_{v,p}(T;0)\ge T^{p-1/2-\alpha}$ for all sufficiently large $T$.
\end{enumerate}
\end{stmtP}
\leanline{\texttt{hs2}, \texttt{dimer\_\allowbreak{}second\_\allowbreak{}moment}, \texttt{integral\_\allowbreak{}hs2\_\allowbreak{}tmProd\_\allowbreak{}le}, \texttt{dimer\_\allowbreak{}super\_\allowbreak{}annealed}, \texttt{dimer\_\allowbreak{}super\_\allowbreak{}ae} \,(\texttt{DimerSecond.\allowbreak{}lean}, \texttt{DimerSuper.\allowbreak{}lean})}
\renewcommand{\stmtname}{\Cref{thm:dimer}}
\begin{stmtP}[Delocalization for the dimer sector law]
Let $0<2|h|<J$, $\lambda:=2|h|/J\in(0,1)$, and let $c_{\lambda,p}$ be as in \eqref{eq:clp}.  For even $N$
and $x\in\{1,\dots,N\}$ let $\Psi^x_N=\cdr_x^\dagger\Omdim$ be the single-fermion dimer state of
\cref{def:dimerstate} (a product state), $\Psi^x_N(t):=\e^{-\ii tH}\Psi^x_N$, and for $p>0$, $T>0$
\begin{equation}\tag{\ref{eq:LGTmoments}}
  D^{(N)}_p(t;x):=\Bigl\langle\Psi^x_N(t),\sum_{y=1}^N|y-x|^p\,n_y\,\Psi^x_N(t)\Bigr\rangle,\qquad
  \Mom^{(N)}_p(T;x):=\int_0^\infty\frac{\dd t}{T}\,\e^{-t/T}D^{(N)}_p(t;x).
\end{equation}
\begin{enumerate}[label=(\alph*)]
\item \emph{(Every sector, every volume.)}  For every even $N$, every $q\in\supp\mudim$, every $x$, every
$p>0$ and every $T$ with $1\le JT\le\min\{x-1,N-x\}$,
\begin{equation}\tag{\ref{eq:dimerLGTa}}
  \int_0^\infty\frac{\dd t}{T}\,\e^{-t/T}\sum_{y=1}^N|y-x|^p\,\bigl|\langle\delta_y,\e^{-\ii t\hq}\delta_x\rangle\bigr|^2
  \ge c_{\lambda,p}\,(JT)^{p-1}.
\end{equation}
Consequently $\Mom^{(N)}_p(T;x)\ge c_{\lambda,p}(JT)^{p-1}$ under the same conditions.
\item \emph{(Thermodynamic limit.)}  Let $x_N\in\{1,\dots,N\}$ with $x_N\to\infty$, $N-x_N\to\infty$, and
$x_N\equiv r+1\pmod2$ for a fixed $r\in\{0,1\}$.  Then for $0<p\le2$, $t\in\R$ and $T>0$,
\begin{align*}
  \lim_{N\to\infty}D^{(N)}_p(t;x_N)&=\E\sum_{y\in\Z}|y|^p\bigl|\langle\delta_y,\e^{-\ii JtH^{(r)}}\delta_0\rangle\bigr|^2,\\
  \lim_{N\to\infty}\Mom^{(N)}_p(T;x_N)&=\E\,\Mom^\Z_{v^{(r)},p}(JT;0)=:\Mom^{\infty,r}_p(T).
\end{align*}
\item \emph{(Superdiffusion.)}  Let $0<p\le2$ and $r\in\{0,1\}$.  Then
$\Mom^{\infty,r}_p(T)\ge c_{\lambda,p}(JT)^{p-1}$ for $JT\ge1$; there are $c',T_0>0$, depending only on $\lambda$ and $p$, with
$\Mom^{\infty,r}_p(T)\ge c'(JT)^{p-1/2}$ for $JT\ge T_0$; and
\[
  \liminf_{T\to\infty}\frac{\log\Mom^\Z_{v^{(r)},p}(T;0)}{\log T}\ge p-\frac12\quad\text{almost surely},\qquad
  \liminf_{T\to\infty}\frac{\log\Mom^{\infty,r}_p(T)}{\log T}\ge p-\frac12 .
\]
In particular, $\Mom^{\infty,r}_2(T)\ge c'(JT)^{3/2}$ for $JT\ge T_0$.
\item \emph{(Failure of uniform localization.)}  For every $p>1/2$,
\begin{equation}\tag{\ref{eq:dimerLGTd}}
  \sup_{N\ \mathrm{even}}\ \sup_{1\le x\le N}\ \sup_{t\in\R}\ \Edim\sum_{y=1}^N|y-x|^p\,
  \bigl|\langle\delta_y,\e^{-\ii t\hq}\delta_x\rangle\bigr|^2=\infty .
\end{equation}
Hence there is no $F:\N_0\to[0,1]$ with $\sum_{d\ge0}(d+1)^pF(d)<\infty$ for some $p>1/2$ such that
$\Edim Q^{[1,N]}_q(x,y)\le F(|x-y|)$ for all even $N$ and all $x,y\in\{1,\dots,N\}$.  In particular,
Hypothesis~\ref{hyp:EC} with the i.i.d.\ law replaced by $\mudim$ fails.  For $p>1$ the supremum is
infinite sector by sector: for every choice of $q^{(N)}\in\supp\mudim$,
$\sup_{N,x,t}\sum_y|y-x|^p|\langle\delta_y,\e^{-\ii th_{q^{(N)}}}\delta_x\rangle|^2=\infty$.
\end{enumerate}
\end{stmtP}
\leanline{\texttt{dimer\_\allowbreak{}sector\_\allowbreak{}moment\_\allowbreak{}lower}, \texttt{dimerMN\_\allowbreak{}lower}, \texttt{dimer\_\allowbreak{}LGT\_\allowbreak{}moment\_\allowbreak{}lower}, \texttt{dimerDLGT\_\allowbreak{}eq}, \texttt{dimer\_\allowbreak{}TL\_\allowbreak{}D}, \texttt{dimer\_\allowbreak{}TL\_\allowbreak{}M}, \texttt{dimer\_\allowbreak{}LGT\_\allowbreak{}TL\_\allowbreak{}D}, \texttt{dimer\_\allowbreak{}LGT\_\allowbreak{}TL\_\allowbreak{}M}, \texttt{dimerMInf\_\allowbreak{}lower}, \texttt{dimerMInf\_\allowbreak{}super}, \texttt{dimerMInf\_\allowbreak{}exponent}, \texttt{dimer\_\allowbreak{}sup\_\allowbreak{}moment\_\allowbreak{}unbounded}, \texttt{dimerD\_\allowbreak{}unbounded\_\allowbreak{}half}, \texttt{dimer\_\allowbreak{}LGT\_\allowbreak{}unbounded\_\allowbreak{}half}, \texttt{not\_\allowbreak{}EC\_\allowbreak{}dimer\_\allowbreak{}half} \,(\texttt{DimerLGT.\allowbreak{}lean}, \texttt{DimerInfVol.\allowbreak{}lean}, \texttt{DimerState.\allowbreak{}lean}, \texttt{DimerSuper.\allowbreak{}lean})}
\renewcommand{\stmtname}{\Cref{prop:dimerimbalance}}
\begin{stmtP}[The imbalance of the dimer CDW state vanishes]
Let $h\ne0$, $\lambda:=2h/J$, $N=2M$, $\Gamma:=\diag((-1)^x)_{x=1}^N$, $\Psi(t):=\e^{-\ii tH}\Psidim(\ncdw)$,
and $I_N(t):=\frac2N\sum_{x=1}^N(-1)^x\langle\Psi(t),n_x\Psi(t)\rangle$.
\begin{enumerate}[label=(\roman*)]
\item \emph{(Identity.)}  For every $q\in\supp\mudim$ and every real normalized eigenvector $\varphi$ of
$\hq/J=-\Delta_N+\lambda\diag(q)$ with eigenvalue $E$,
\begin{equation}\tag{\ref{eq:dimerimbid}}
  (\lambda^2-E^2)\sum_{x=1}^N(-1)^x\varphi(x)^2=\varphi(N)^2-\varphi(1)^2 .
\end{equation}
\item The limit $\bar I_N:=\lim_{T\to\infty}T^{-1}\int_0^TI_N(t)\,\dd t$ exists, and
\[
  \bar I_N=\Edim\frac1N\sum_\alpha a_\alpha^2,\qquad a_\alpha:=\langle\varphi_\alpha,\Gamma\varphi_\alpha\rangle,
\]
for a real orthonormal eigenbasis $(\varphi_\alpha)$ of $\hq$.
\item For every integer $L\ge1$,
\begin{equation}\tag{\ref{eq:dimerimbrate}}
  \bar I_N\le\frac{8L^2(2|\lambda|+1)^{2L}}{N}+\frac4L+\frac4N .
\end{equation}
In particular $\lim_{N\to\infty}\bar I_N=0$.
\end{enumerate}
\end{stmtP}
\leanline{\texttt{dimer\_\allowbreak{}stagWeight\_\allowbreak{}identity}, \texttt{dimerImbalance\_\allowbreak{}timeAvg}, \texttt{card\_\allowbreak{}eig\_\allowbreak{}near\_\allowbreak{}le}, \texttt{dimerImbInf\_\allowbreak{}le}, \texttt{dimerImbInf\_\allowbreak{}tendsto\_\allowbreak{}zero} \,(\texttt{DimerImbalance.\allowbreak{}lean}, \texttt{DimerImbalanceLim.\allowbreak{}lean})}
\renewcommand{\stmtname}{\Cref{lem:fewapprox}}
\begin{stmtP}[Few eigenvalues near a given energy]
Let $N,L\ge1$, $v:[1,N]\to\R$, $E\in\R$, $K:=\max_y|v(y)-E|$ and $H:=-\Delta_N+\diag(v)$.  If $\varepsilon\ge0$
and $\varepsilon L(K+1)^L<1$, then $H$ has at most $2(N/L+1)$ eigenvalues in $[E-\varepsilon,E+\varepsilon]$.
\end{stmtP}
\leanline{\texttt{dimer\_\allowbreak{}stagWeight\_\allowbreak{}identity}, \texttt{dimerImbalance\_\allowbreak{}timeAvg}, \texttt{card\_\allowbreak{}eig\_\allowbreak{}near\_\allowbreak{}le}, \texttt{dimerImbInf\_\allowbreak{}le}, \texttt{dimerImbInf\_\allowbreak{}tendsto\_\allowbreak{}zero} \,(\texttt{DimerImbalance.\allowbreak{}lean}, \texttt{DimerImbalanceLim.\allowbreak{}lean});\\ \texttt{dosFun}, \texttt{dosFin}, \texttt{tendsto\_\allowbreak{}dosFin}, \texttt{dosFun\_\allowbreak{}le\_\allowbreak{}of\_\allowbreak{}near}, \texttt{exists\_\allowbreak{}dosMeasure}, \texttt{dosMeasure}, \texttt{isProbabilityMeasure\_\allowbreak{}dosMeasure}, \texttt{dosMeasure\_\allowbreak{}compl\_\allowbreak{}Icc}, \texttt{dosMeasure\_\allowbreak{}singleton} \,(\texttt{DimerDOS.\allowbreak{}lean})}
\subsection*{\S\,\ref{sec:2d}\quad Two dimensions}
\addcontentsline{toc}{subsection}{\S\,\ref{sec:2d}\quad Two dimensions}
\renewcommand{\stmtname}{\Cref{lem:2dgauge}}
\begin{stmtP}[Fluxes and gauge orbits]
\leavevmode
\begin{enumerate}[label=(\roman*)]
\item For $s,s'\in\{\pm1\}^E$: $\varphi_s=\varphi_{s'}$ if and only if $s'=\theta\cdot s$ for some
$\theta\in\TGauge$, and $\theta$ is then unique up to a global sign.
\item Let $T\subset E$ be a spanning tree.  For $b\in E\setminus T$ let $\gamma_b$ be the edge set
consisting of $b$ and the path in $T$ joining the endpoints of $b$.  Then
$\{\gamma_b:b\in E\setminus T\}$ is a basis of $\TCyc(\TLat)$; in particular
$\dim_{\mathbb F_2}\TCyc(\TLat)=\beta$ and $\abs{\TFlux(\TLat)}=2^\beta$.
\item The map $s\mapsto\varphi_s$ is onto $\TFlux(\TLat)$, and each fibre
$O_\varphi:=\{s:\varphi_s=\varphi\}$ is a $\TGauge$-orbit with $2^{\abs V-1}$ elements.
\end{enumerate}
\end{stmtP}
\leanline{\texttt{FinGraph.\allowbreak{}flux\_\allowbreak{}eq\_\allowbreak{}iff\_\allowbreak{}gauge}, \texttt{FinGraph.\allowbreak{}gaugeAct\_\allowbreak{}eq\_\allowbreak{}self\_\allowbreak{}iff}, \texttt{FinGraph.\allowbreak{}card\_\allowbreak{}gaugeOrbit}, \texttt{FinGraph.\allowbreak{}card\_\allowbreak{}fluxClasses}, \texttt{FinGraph.\allowbreak{}wilsonLoop\_\allowbreak{}comm\_\allowbreak{}ham}, \texttt{FinGraph.\allowbreak{}wilsonLoop\_\allowbreak{}comm\_\allowbreak{}gaussOp}, \texttt{FinGraph.\allowbreak{}gaussOp\_\allowbreak{}prod} \,(\texttt{GraphFlux.\allowbreak{}lean}, \texttt{GraphWilson.\allowbreak{}lean})}
\renewcommand{\stmtname}{\Cref{thm:2dmap}}
\begin{stmtP}[Exact mapping on graphs; planar boxes and tori]
Let $\TLat=(V,E)$ be a finite connected graph without loops.
\begin{enumerate}[label=(\alph*)]
\item The $G_i$ ($i\in V$) and $W_\gamma$ ($\gamma\in\TCyc(\TLat)$) are self-adjoint unitaries; they commute
with each other, with $N_f$ and with $H_\TLat$; and
\begin{equation}\tag{\ref{eq:2dparity}}
  \prod_{i\in V}G_i=(-1)^{N_f}.
\end{equation}
\item For $q\in\{\pm1\}^V$ and $\varphi\in\TFlux(\TLat)$ let $P_{q,\varphi}$ be the joint spectral projection
onto $\{G_i=q_i\ \forall i,\ W_\gamma=\varphi(\gamma)\ \forall\gamma\}$, and let $\Tpi_q:=\prod_iq_i$.  Then
$\sum_{q,\varphi}P_{q,\varphi}=\one$, $\operatorname{rank}P_{q,\varphi}=2^{\abs V-1}$ for \emph{every}
$(q,\varphi)$, and $(-1)^{N_f}P_{q,\varphi}=\Tpi_qP_{q,\varphi}$.  For every $s$ with $\varphi_s=\varphi$ there
is a unitary $V_{q,s}:\cF_V^{(\Tpi_q)}\to\operatorname{Ran}P_{q,\varphi}$ which maps each $\ket n$ to a joint
eigenvector of the $n_i$ with the same eigenvalues, and such that, on $\cF_V^{(\Tpi_q)}$,
\begin{equation}\tag{\ref{eq:2ddict}}
  V_{q,s}^*\,\sigma^z_bf_i^\dagger f_j\,V_{q,s}=s_b\,f_i^\dagger f_j\quad(\partial b=\{i,j\}),\qquad
  V_{q,s}^*A_iV_{q,s}=q_i(-1)^{n_i},
\end{equation}
and consequently
\begin{equation}\tag{\ref{eq:2dsector}}
  V_{q,s}^*\,H_\TLat\,V_{q,s}=\Bigl(\dGamma(h_{q,s})-h\sum_{i\in V}q_i\Bigr)\Big|_{\cF_V^{(\Tpi_q)}} .
\end{equation}
For $\theta\in\TGauge$, $V_{q,\theta\cdot s}=q(\theta)\,V_{q,s}\,\Gamma(D_\theta)$ with
$q(\theta):=\prod_{i:\theta_i=-1}q_i$.
\item \emph{Planar box.}  Let $\TLat$ be the box $\{1,\dots,L_1\}\times\{1,\dots,L_2\}$ with nearest-neighbour
edges, and let the plaquettes $p$ be its unit squares, with boundaries $\partial p$ and $W_p:=W_{\partial p}$.
Then $\beta=(L_1-1)(L_2-1)$, and $\{\partial p\}_p$ is a basis of $\TCyc(\TLat)$.  Hence the sectors are
labelled by $(q,(W_p)_p)\in\{\pm1\}^V\times\{\pm1\}^{\#\text{plaquettes}}$, all $2^{\abs V+\beta}$ labels
occur, and in the sector $(q,(W_p))$ the matter consists of free fermions with the binary potential $2hq_i$
and the binary $\pi$-flux $W_p=\prod_{b\in\partial p}s_b\in\{\pm1\}$ through $p$, restricted to the parity
sector $(-1)^{N_f}=\prod_iq_i$.
\item \emph{Torus.}  Let $\TLat$ be the torus $(\Z/L_1\Z)\times(\Z/L_2\Z)$, $L_1,L_2\ge3$, with
nearest-neighbour edges, plaquettes $p$ as in (c), and let $\gamma_1$ (resp.\ $\gamma_2$) be the set of
horizontal edges of the row $\{y=0\}$ (resp.\ the vertical edges of the column $\{x=0\}$).  Then
$\beta=L_1L_2+1$, $\prod_pW_p=\one$, and $\{\partial p:p\neq p_0\}\cup\{\gamma_1,\gamma_2\}$ is a basis of
$\TCyc(\TLat)$ for any fixed plaquette $p_0$.  Hence the sectors are labelled by $q$, the plaquette fluxes
$(W_p)_p$ subject to $\prod_pW_p=1$, and the two non-contractible Wilson loops $W_{\gamma_1},W_{\gamma_2}$,
which are independent labels; all $2^{\abs V+\beta}$ labels occur.  In $h_{q,s}$ the non-contractible loops are the
products of the hopping signs along $\gamma_1,\gamma_2$; when all $W_p=1$ they act as $\Z_2$ twists of the boundary
conditions in the two directions.
\end{enumerate}
\end{stmtP}
\leanline{\texttt{FinGraph.\allowbreak{}sector\_\allowbreak{}average\_\allowbreak{}matter}, \texttt{FinGraph.\allowbreak{}gaussOp\_\allowbreak{}comm\_\allowbreak{}ham}, \texttt{FinGraph.\allowbreak{}dGamma\_\allowbreak{}hq} \,(\texttt{Graph.\allowbreak{}lean});\\ \texttt{FinGraph.\allowbreak{}flux\_\allowbreak{}eq\_\allowbreak{}iff\_\allowbreak{}gauge}, \texttt{FinGraph.\allowbreak{}gaugeAct\_\allowbreak{}eq\_\allowbreak{}self\_\allowbreak{}iff}, \texttt{FinGraph.\allowbreak{}card\_\allowbreak{}gaugeOrbit}, \texttt{FinGraph.\allowbreak{}card\_\allowbreak{}fluxClasses}, \texttt{FinGraph.\allowbreak{}wilsonLoop\_\allowbreak{}comm\_\allowbreak{}ham}, \texttt{FinGraph.\allowbreak{}wilsonLoop\_\allowbreak{}comm\_\allowbreak{}gaussOp}, \texttt{FinGraph.\allowbreak{}gaussOp\_\allowbreak{}prod} \,(\texttt{GraphFlux.\allowbreak{}lean}, \texttt{GraphWilson.\allowbreak{}lean});\\ \texttt{FinGraph.\allowbreak{}fluxVec\_\allowbreak{}gauss}, \texttt{FinGraph.\allowbreak{}fluxVec\_\allowbreak{}wilson}, \texttt{FinGraph.\allowbreak{}fluxVec\_\allowbreak{}occ}, \texttt{FinGraph.\allowbreak{}fluxVec\_\allowbreak{}gaugeAct}, \texttt{FinGraph.\allowbreak{}fluxVec\_\allowbreak{}norm}, \texttt{FinGraph.\allowbreak{}fluxVec\_\allowbreak{}orthogonal}, \texttt{FinGraph.\allowbreak{}card\_\allowbreak{}fluxLabels}, \texttt{FinGraph.\allowbreak{}ham\_\allowbreak{}fluxVec} \,(\texttt{TwoDMap.\allowbreak{}lean});\\ \texttt{FinGraph.\allowbreak{}IsCycleBasis}, \texttt{FinGraph.\allowbreak{}gauge\_\allowbreak{}iff\_\allowbreak{}flux\_\allowbreak{}basis}, \texttt{FinGraph.\allowbreak{}exists\_\allowbreak{}flux\_\allowbreak{}basis}, \texttt{FinGraph.\allowbreak{}card\_\allowbreak{}cycleBasis}, \texttt{FinGraph.\allowbreak{}box\_\allowbreak{}connected}, \texttt{FinGraph.\allowbreak{}box\_\allowbreak{}beta}, \texttt{FinGraph.\allowbreak{}box\_\allowbreak{}plaq\_\allowbreak{}basis}, \texttt{FinGraph.\allowbreak{}box\_\allowbreak{}card\_\allowbreak{}fluxClasses}, \texttt{FinGraph.\allowbreak{}torus\_\allowbreak{}connected}, \texttt{FinGraph.\allowbreak{}torus\_\allowbreak{}beta}, \texttt{FinGraph.\allowbreak{}torus\_\allowbreak{}sum\_\allowbreak{}plaq}, \texttt{FinGraph.\allowbreak{}torus\_\allowbreak{}prod\_\allowbreak{}wilson}, \texttt{FinGraph.\allowbreak{}torus\_\allowbreak{}basis}, \texttt{FinGraph.\allowbreak{}torus\_\allowbreak{}card\_\allowbreak{}fluxClasses} \,(\texttt{TwoDLattices.\allowbreak{}lean})}
\renewcommand{\stmtname}{\Cref{prop:2dlaws}}
\begin{stmtP}[Sector laws]
Let $\TLat$ be as in \cref{thm:2dmap}.
\begin{enumerate}[label=(\alph*)]
\item \emph{$\sigma^z$-product links: the Anderson--Bernoulli model.}  Let $\Psi_0=\psi\otimes\ket{s}_z$ with
$s\in\{\pm1\}^E$ and $\psi\in\cF_V^{(\Tpi)}$ of definite parity $\Tpi$.  Then $P_{q,\varphi}\Psi_0=0$ unless
$\varphi=\varphi_s$ and $\Tpi_q=\Tpi$, and for these sectors $P_{q,\varphi_s}\Psi_0=2^{-(\abs V-1)/2}\,V_{q,s}\psi$.  So the
flux is deterministic, the sector law of the charges is uniform on the parity hyperplane $\{q:\prod_iq_i=\Tpi\}$,
every sector carries the matter state $\psi$, and for every $O$ commuting with all $G_i,W_\gamma$
\begin{equation}\tag{\ref{eq:2dlawz}}
  \langle\Psi(t),O\Psi(t)\rangle=2^{1-\abs V}\sum_{q:\ \prod_iq_i=\Tpi}
  \bigl\langle\Gamma(\e^{-\ii th_{q,s}})\psi,\,O_{q,s}\,\Gamma(\e^{-\ii th_{q,s}})\psi\bigr\rangle .
\end{equation}
The law of $q$ has density at most $2$ with respect to the i.i.d.\ Bernoulli$(1/2)$ law.  For the
$\sigma^z$-polarized state $\ket{\uparrow\cdots\uparrow}_z$, i.e.\ $s=+:=(1,\dots,1)$, all fluxes vanish and the
one-body operators are the Anderson--Bernoulli operators $h_{q,+}=-J\Delta+2h\diag(q)$, with $\Delta$ the adjacency
operator of $\TLat$.
\item \emph{$\sigma^x$-polarized links: random $\pi$-flux.}  Let $\Psi_0=\ket{n^0}\otimes\ket{+\cdots+}$, where
$\sigma^x_b\ket{+\cdots+}=\ket{+\cdots+}$ for all $b$.  Let $q^0_i:=(-1)^{n^0_i}$.  Then $P_{q,\varphi}\Psi_0=0$
unless $q=q^0$, and for every $\varphi\in\TFlux(\TLat)$ and every $s\in O_\varphi$,
$P_{q^0,\varphi}\Psi_0=2^{-\beta/2}\,V_{q^0,s}\ket{n^0}$.  So the charges are deterministic, the flux
configuration is uniformly distributed on $\TFlux(\TLat)$ (on a box: i.i.d.\ uniform plaquette fluxes; on a
torus: $(W_p)$ uniform on $\{\prod_pW_p=1\}$, and $W_{\gamma_1},W_{\gamma_2}$ independent and uniform), every sector
carries the matter state $\ket{n^0}$, and
\begin{equation}\tag{\ref{eq:2dlawx}}
  \langle\Psi(t),O\Psi(t)\rangle=2^{-\beta}\sum_{\varphi\in\TFlux}
  \bigl\langle\Gamma(\e^{-\ii th_{q^0,s^\varphi}})\ket{n^0},\,O_{q^0,s^\varphi}\,\Gamma(\e^{-\ii th_{q^0,s^\varphi}})\ket{n^0}\bigr\rangle ,
  \qquad s^\varphi\in O_\varphi .
\end{equation}
\item \emph{Checkerboard.}  In (b) let $\TLat$ be bipartite with sublattice signs $\epsilon_i\in\{\pm1\}$
($\epsilon_i=-\epsilon_j$ whenever $\partial b=\{i,j\}$), and $n^0_i=\frac12(1+\epsilon_i)$.  Let
$\Gamma_\epsilon:=\diag(\epsilon)$.  Then $h_{q^0,s}=K_s-2h\Gamma_\epsilon$ with
$K_s\Gamma_\epsilon=-\Gamma_\epsilon K_s$, and for every $s$:
\begin{enumerate}[label=(\roman*)]
\item $h_{q^0,s}^2=K_s^2+4h^2$;
\item $\sigma(h_{q^0,s})\subset(-\infty,-2\abs h]\cup[2\abs h,\infty)$, and for every Borel set $B\subset[0,\infty)$,
$\one_B(\abs{h_{q^0,s}})=\one_{\omega^{-1}(B)}(\abs{K_s})$ with $\omega(e):=\sqrt{e^2+4h^2}$ ($e\ge0$).  In particular
$\one_{\{2\abs h\}}(\abs{h_{q^0,s}})$ is the projection onto $\ker K_s$, and
$\one_{[2\abs h,\sqrt{4h^2+\delta^2}]}(\abs{h_{q^0,s}})=\one_{[0,\delta]}(\abs{K_s})$ for $\delta\ge0$: the
chiral point $E=0$ of the random-flux operator $K_s$ is mapped to the inner band edges $\abs E=2\abs h$;
\item with $R_s:=(K_s^2+4h^2)^{1/2}$, $\e^{-\ii th_{q^0,s}}=\cos(tR_s)-\ii\,h_{q^0,s}R_s^{-1}\sin(tR_s)$;
\item if $\TLat$ is simple (no multiple edges), $(K_s^2)_{ii}=J^2\deg(i)$.  (For parallel edges of opposite signs
the hopping amplitudes cancel, and this fails.)  If $\TLat$ is a box in $\Z^2$, $K_s^2$ connects only sites of the same
sublattice at distance $2$: $(K_s^2)_{ik}=J^2s_{ij}s_{jk}$ if $k=i\pm2e_\mu$ (with midpoint $j$), and
$(K_s^2)_{ik}=J^2s_{ij_1}s_{j_1k}(1+W_p)$ if $i,k$ are opposite corners of the plaquette $p$ with the other
corners $j_1,j_2$; here $s_{ij}$ is the sign of the edge joining $i$ and $j$.  Thus the diagonal amplitudes of
$K_s^2$ vanish across $\pi$-flux plaquettes.
\end{enumerate}
\end{enumerate}
\end{stmtP}
\leanline{\texttt{FinGraph.\allowbreak{}sector\_\allowbreak{}average\_\allowbreak{}flux}, \texttt{FinGraph.\allowbreak{}sector\_\allowbreak{}average\_\allowbreak{}xpol}, \texttt{FinGraph.\allowbreak{}xpol\_\allowbreak{}summand\_\allowbreak{}gauge} \,(\texttt{GraphFlux.\allowbreak{}lean});\\ \texttt{FinGraph.\allowbreak{}checkerboard\_\allowbreak{}sq}, \texttt{FinGraph.\allowbreak{}checkerboard\_\allowbreak{}gap}, \texttt{FinGraph.\allowbreak{}checkerboard\_\allowbreak{}edge\_\allowbreak{}iff}, \texttt{FinGraph.\allowbreak{}Ks\_\allowbreak{}anticomm}, \texttt{FinGraph.\allowbreak{}gauge\_\allowbreak{}hqs}, \texttt{FinGraph.\allowbreak{}hqs\_\allowbreak{}one} \,(\texttt{TwoD.\allowbreak{}lean});\\ \texttt{FinGraph.\allowbreak{}checkerboard\_\allowbreak{}evo\_\allowbreak{}eigvec}, \texttt{FinGraph.\allowbreak{}checkerboard\_\allowbreak{}evo}, \texttt{FinGraph.\allowbreak{}Ks\_\allowbreak{}sq\_\allowbreak{}diag} \,(\texttt{TwoDLaws.\allowbreak{}lean})}
\renewcommand{\stmtname}{\Cref{prop:cluster}}
\begin{stmtP}[Cluster confinement at strong coupling]
Assume $\abs h>2J$ and set $g:=4\abs h-8J>0$.
\begin{enumerate}[label=(\alph*)]
\item $h^0_{q,s}$ leaves $\ell^2(\Tclus)$ invariant for every equal-charge cluster $\Tclus$.
\item For all $t\in\R$, uniformly in $\Lambda$, $q$ and $s$,
\begin{equation}\tag{\ref{eq:clusterbound}}
  \norm{\e^{-\ii th_{q,s}}-\e^{-\ii th^0_{q,s}}}\le\frac{8J}{g}+\frac{16\,\abs t\,J^2}{g} .
\end{equation}
Consequently, for every $x\in\Lambda$, $\norm{\one_{\Lambda\setminus\Tclus_q(x)}\e^{-\ii th_{q,s}}\delta_x}\le8J/g+16\abs tJ^2/g$.
\item Let the law of $q$ have density at most $K$ with respect to the i.i.d.\ Bernoulli$(1/2)$ law on
$\{\pm1\}^\Lambda$ ($K=1$ for i.i.d.\ charges, $K=2$ for the parity hyperplane of \cref{prop:2dlaws}(a)).
There is a numerical constant $c_*>0$ such that for all $x\in\Lambda$, all integers $R\ge4$, $t\in\R$ and $s$,
\begin{equation}\tag{\ref{eq:clusterprob}}
  \E\,\norm{\one_{\Lambda\setminus B_R(x)}\,\e^{-\ii th_{q,s}}\delta_x}^2
  \le\Bigl(\frac{8J}{g}+\frac{16\,\abs t\,J^2}{g}\Bigr)^2+2K\,\e^{-c_*(2R+3)} .
\end{equation}
\end{enumerate}
\end{stmtP}
\leanline{\texttt{FinGraph.\allowbreak{}hqs0\_\allowbreak{}cluster}, \texttt{FinGraph.\allowbreak{}cluster\_\allowbreak{}bound}, \texttt{FinGraph.\allowbreak{}cluster\_\allowbreak{}leak}, \texttt{FinGraph.\allowbreak{}cluster\_\allowbreak{}avg} \,(\texttt{Cluster.\allowbreak{}lean})}
\renewcommand{\stmtname}{\Cref{prop:nocompact}}
\begin{stmtP}[No compactly supported eigenvectors on $\Z^2$]
Let $a_{xy}\in\C\setminus\{0\}$ for every nearest-neighbour pair $\{x,y\}$ of $\Z^2$, let $V:\Z^2\to\C$, and let
$(Hu)(x):=\sum_{y\sim x}a_{xy}u(y)+V(x)u(x)$ for finitely supported $u$.  If $u$ is finitely supported and $Hu=Eu$ for
some $E\in\C$, then $u=0$.  In particular, for no $q$ and no $s$ does $h_{q,s}$ on $\ell^2(\Z^2)$ have a finitely
supported eigenvector.
\end{stmtP}
\leanline{\texttt{no\_\allowbreak{}compact\_\allowbreak{}eigenvector}, \texttt{no\_\allowbreak{}compact\_\allowbreak{}eigenvector\_\allowbreak{}hqs} \,(\texttt{TwoD.\allowbreak{}lean})}
\subsection*{\S\,\ref{sec:u1}\quad Number entropy in \texorpdfstring{$U(1)$}{U(1)} gauge theories: kinematic bounds}
\addcontentsline{toc}{subsection}{\S\,\ref{sec:u1}\quad Number entropy in \texorpdfstring{$U(1)$}{U(1)} gauge theories: kinematic bounds}
\renewcommand{\stmtname}{\Cref{prop:gausscap}}
\begin{stmtP}[Gauss-law caps]
Let $A=[a,b]\subset\{1,\dots,L\}$ be an interval and $\UQ_A:=\sum_{j\in A}G_j$.
\begin{enumerate}[label=(\alph*)]
\item (Summed Gauss law; any links.)  As an operator identity,
\begin{equation}\tag{\ref{eq:summedgauss}}
  N_A=E_b-E_{a-1}+\Uodd_A-\UQ_A .
\end{equation}
\item (Fixed sector; spin-$S$ links.)  Let $\varrho$ be \emph{any} state with $P_Q\varrho P_Q=\varrho$ for some $Q$.  Then the
law of $N_A$ in $\varrho$ is the law of $E_b-E_{a-1}$ shifted by the constant $\Uodd_A-\sum_{j\in A}Q_j$, and
\begin{equation}\tag{\ref{eq:cap}}
  \SN(\varrho;A)\le
  \begin{cases}
    \ln\min\{4S+1,\ \abs A+1\} & (a\ge2),\\
    \ln\min\{2S+1,\ \abs A+1\} & (a=1,\ \text{fixed-flux edge}).
  \end{cases}
\end{equation}
\item (Superpositions of sectors; spin-$S$ links.)  For every state $\varrho$,
\[
  \SN(\varrho;A)\le\ln(4S+1)+\Ush_\varrho(\UQ_A)\quad(a\ge2),\qquad
  \SN(\varrho;A)\le\ln(2S+1)+\Ush_\varrho(\UQ_A)\quad(a=1).
\]
So superposing background-charge sectors adds at most $\Ush_\varrho(\UQ_A)$, the entropy of the total background charge in $A$.
\item (Any dynamics.)  Let $(V(t))_{t\in\R}$ be unitaries commuting with all $G_j$, e.g.\ generated by any
gauge-invariant, possibly time-dependent, interacting or disordered Hamiltonian, and $\varrho(t):=V(t)\varrho V(t)^*$.
If $\varrho$ lies in one sector, then \eqref{eq:cap} holds for $\varrho(t)$ for all $t$; in (c),
$\Ush_{\varrho(t)}(\UQ_A)=\Ush_\varrho(\UQ_A)$ for all $t$.
\end{enumerate}
\end{stmtP}
\leanline{\texttt{energy\_\allowbreak{}conservation\_\allowbreak{}bound}, \texttt{spinPlus\_\allowbreak{}opNorm}, \texttt{max\_\allowbreak{}spinPlus\_\allowbreak{}weight\_\allowbreak{}even}, \texttt{max\_\allowbreak{}spinPlus\_\allowbreak{}weight\_\allowbreak{}odd}, \texttt{cS\_\allowbreak{}bounds}, \texttt{shannon\_\allowbreak{}uniform}, \texttt{qlm\_\allowbreak{}numberEntropy\_\allowbreak{}attained} \,(\texttt{U1Remarks.\allowbreak{}lean});\\ \texttt{countInt\_\allowbreak{}gauss}, \texttt{qlm\_\allowbreak{}numberEntropy\_\allowbreak{}le}, \texttt{qlm\_\allowbreak{}numberEntropy\_\allowbreak{}le\_\allowbreak{}edge}, \texttt{qlm\_\allowbreak{}numberEntropy\_\allowbreak{}le\_\allowbreak{}super}, \texttt{qlm\_\allowbreak{}numberEntropy\_\allowbreak{}le\_\allowbreak{}super\_\allowbreak{}edge}, \texttt{qlm\_\allowbreak{}sector\_\allowbreak{}preserved}, \texttt{bgEntropy\_\allowbreak{}invariant} \,(\texttt{QLM.\allowbreak{}lean}, \texttt{QLMSuper.\allowbreak{}lean})}
\renewcommand{\stmtname}{\Cref{prop:wilson}}
\begin{stmtP}[Wilson links: cut-averaged number entropy]
Consider Wilson links and $H$ as in \eqref{eq:HU1}.  Let $\Psi_0\in\operatorname{Ran}P_Q$ be a unit vector,
$\Psi(t):=\e^{-\ii tH}\Psi_0$ and $\varepsilon_0:=\langle\Psi_0,H\Psi_0\rangle/L$.  Then
$\varepsilon_0+w+\abs m\ge0$, and for all $t\in\R$ and all $L\ge1$:
\begin{enumerate}[label=(\alph*)]
\item $\SN(\Psi(t);[1,\ell])=\Ush_{\Psi(t)}(E_\ell)$ for $\ell=1,\dots,L$;
\item $\displaystyle\frac1L\sum_{\ell=1}^L\langle\Psi(t),(E_\ell+\vartheta)^2\Psi(t)\rangle\le\frac2{g^2}\bigl(\varepsilon_0+w+\abs m\bigr)$;
\item $\displaystyle\frac1L\sum_{\ell=1}^L\SN(\Psi(t);[1,\ell])\le\frac12\ln\Bigl(2\pi e\Bigl(\frac2{g^2}(\varepsilon_0+w+\abs m)+\frac1{12}\Bigr)\Bigr)$;
\item for every $\delta\in(0,1]$, $\SN(\Psi(t);[1,\ell])\le\frac12\ln\bigl(2\pi e\bigl(\frac{2(\varepsilon_0+w+\abs m)}{\delta g^2}+\frac1{12}\bigr)\bigr)$
for all but at most $\delta L$ values of $\ell$;
\item for every $\ell$, $\SN(\Psi(t);[1,\ell])\le\frac12\ln\bigl(2\pi e\bigl(\frac{2L(\varepsilon_0+w+\abs m)}{g^2}+\frac1{12}\bigr)\bigr)$.
\end{enumerate}
\end{stmtP}
\leanline{\texttt{wField\_\allowbreak{}hop}, \texttt{wField\_\allowbreak{}eq\_\allowbreak{}count}, \texttt{wilson\_\allowbreak{}field\_\allowbreak{}bound}, \texttt{wilson\_\allowbreak{}numberEntropy\_\allowbreak{}avg}, \texttt{wilson\_\allowbreak{}numberEntropy\_\allowbreak{}most}, \texttt{wilson\_\allowbreak{}numberEntropy\_\allowbreak{}each} \,(\texttt{Wilson.\allowbreak{}lean});\\ \texttt{wStateS}, \texttt{energyDensityS}, \texttt{lawEntropy}, \texttt{wilson\_\allowbreak{}field\_\allowbreak{}bound\_\allowbreak{}super}, \texttt{bgEntropyS\_\allowbreak{}wStateS}, \texttt{wilson\_\allowbreak{}numberEntropy\_\allowbreak{}le\_\allowbreak{}super}, \texttt{lawEntropy\_\allowbreak{}le\_\allowbreak{}massey}, \texttt{lawEntropy\_\allowbreak{}le\_\allowbreak{}log\_\allowbreak{}card}, \texttt{wilson\_\allowbreak{}numberEntropy\_\allowbreak{}avg\_\allowbreak{}super}, \texttt{bgEntropyS\_\allowbreak{}bare\_\allowbreak{}vacuum}, \texttt{bgEntropyS\_\allowbreak{}le\_\allowbreak{}log\_\allowbreak{}three} \,(\texttt{WilsonSuper.\allowbreak{}lean})}
\renewcommand{\stmtname}{\Cref{prop:brenes}}
\begin{stmtP}[The states of [19]]
Let $\ket{\rm vac}$ be the bare vacuum ($\rho_j\ket{\rm vac}=0$ for all $j$), $e_0\in\Z$, and
\[
  \Psi_0:=\ket{\rm vac}\otimes\bigotimes_{j=1}^L\tfrac1{\sqrt3}\bigl(\ket{-1}+\ket0+\ket1\bigr)_{E_j}
  =3^{-L/2}\sum_{\bar e\in\{-1,0,1\}^L}\ket{\rm vac}\otimes\ket{\bar e}.
\]
\begin{enumerate}[label=(\alph*)]
\item $\ket{\rm vac}\otimes\ket{\bar e}$ lies in the sector with background charges $Q_j=\bar e_j-\bar e_{j-1}$
($\bar e_0:=e_0$), and $\bar e\mapsto Q$ is injective.  So every gauge-invariant expectation value in $\Psi(t)$, at any
number of times, is the uniform average over the \emph{background fields} $\bar e$, which are i.i.d.\ uniform on
$\{-1,0,1\}$, of the corresponding quantity in the sector of $\bar e$.  The background charges are not independent.
\item In the sector of $\bar e$, $E_j=\bar e_j+F_j$ with $F_j:=\sum_{k\le j}\rho_k$, and
\[
  H=H_{\rm cl}+g^2\sum_{k=1}^LA_k\rho_k+\frac{g^2}2\sum_{j=1}^L(\bar e_j+\vartheta)^2,\qquad
  A_k:=\sum_{j\ge k}(\bar e_j+\vartheta),
\]
where $H_{\rm cl}$ is \eqref{eq:HU1} in the sector $Q=0$ with $e_0=\vartheta=0$.  So the charges move in the clean
model plus the random-walk potential $g^2A_k$, whose steps $g^2(\bar e_k+\vartheta)$ are i.i.d.  Creating a meson
from $\ket{\rm vac}$ on the bond $(j,j+1)$ changes the energy by $2m+\frac{g^2}2\bigl(1+2(-1)^j(\bar e_j+\vartheta)\bigr)$:
a tightly bound meson sees an i.i.d.\ potential.
\item $\UQ_{[1,\ell]}=\bar e_\ell-e_0$ on $\Psi_0$, so $\Ush_{\Psi_0}(\UQ_{[1,\ell]})=\ln3$ for every $\ell$, and for all
$t\in\R$ and $L\ge1$
\[
  \frac1L\sum_{\ell=1}^L\SN(\Psi(t);[1,\ell])\le\ln3+\frac12\ln\Bigl(2\pi e\Bigl(\frac2{g^2}\bigl(\varepsilon_0+w+\abs m\bigr)
  +\frac1{12}\Bigr)\Bigr),
\]
where $\varepsilon_0=\frac{g^2}2\bigl(\frac23+\vartheta^2\bigr)-m\lceil L/2\rceil/L$ is the energy density of $\Psi_0$.
\item For all $t\in\R$,
$\bigl|\frac1L\sum_{j=1}^L\langle(E_j+\vartheta)^2\rangle_t-\bigl(\frac23+\vartheta^2\bigr)\bigr|\le4(w+\abs m)/g^2$.
\end{enumerate}
\end{stmtP}
\leanline{\texttt{fieldOf}, \texttt{brenesQ}, \texttt{bareVac}, \texttt{brenesInit}, \texttt{brenesPot}, \texttt{wField\_\allowbreak{}brenes}, \texttt{brenesPot\_\allowbreak{}succ}, \texttt{schwinger\_\allowbreak{}brenes}, \texttt{schwinger\_\allowbreak{}meson}, \texttt{bgSum\_\allowbreak{}brenes}, \texttt{bgEntropyS\_\allowbreak{}brenes}, \texttt{energyDensityS\_\allowbreak{}brenes}, \texttt{brenes\_\allowbreak{}numberEntropy\_\allowbreak{}avg}, \texttt{brenes\_\allowbreak{}field\_\allowbreak{}memory} \,(\texttt{BrenesU1.\allowbreak{}lean})}


\end{document}
